% 中文翻译；英文底本行 10089--11033。来源对齐状态与疑难核验见 audits/expose-06.json。
\sgaStableTarget{sga6.e06.expose}\section*{第 VI 讲. 射影丛的 \(K^\bullet\)：计算及其推论}
\begin{center}
Pierre Berthelot
\end{center}

\sgaStableTarget{sga6.e06.sec.1}\section*{1. 射影丛的 \(K^\bullet\) 的计算：有限型局部自由层的情形}

本节中，\(K^\bullet(X)\) 表示概形 \(X\) 上有限型局部自由层的 Grothendieck 群。我们的目的是证明下述结构定理。

\sgaStableTarget{sga6.e06.thm.1-1}\textbf{定理 1.1.} 设 \(S\) 是拟紧概形，\(E\) 是秩为 \(r+1\) 的局部自由 \(\mathcal O_S\)-模，并设
\[
  X=\mathbf P(E)
\]
为相应的射影丛。结构态射 \(f:X\to S\) 赋予 \(K^\bullet(X)\) 一个 \(K^\bullet(S)\)-代数结构。若 \(\xi\in K^\bullet(X)\) 表示 \(\mathcal O_X(1)\) 的类，则 \(K^\bullet(X)\) 是以
\[
  1,\xi,\ldots,\xi^r .
\]
为基的自由 \(K^\bullet(S)\)-模。

\subsection*{证明的安排}
证明分为三个部分。

首先证明，\(\xi\) 的幂作为 \(K^\bullet(S)\)-模生成 \(K^\bullet(X)\)。为此，只需建立下述事实。若 \(F\) 是 \(X\) 上的有限型局部自由层，则可以找到一个有限表现、且在 \(\mathcal O_S\) 上平坦的分次 \(\Sym(E)\)-模 \(M\)，使得 \(F=\Proj_0(M)\)（\(\Sym(E)\) 是 \(E\) 的对称代数；记号 \(\Proj_0\) 参见 EGA II 8.12.1）。这样的 \(M\) 有一个有限分次消解
\[
  0\longrightarrow L_{r+1}\longrightarrow \cdots \longrightarrow L_0\longrightarrow M\longrightarrow 0
\]
其中各项都是有限型局部自由分次 \(\Sym(E)\)-模。最后，每个 \(L_i\) 有一个有限滤过，其商形如
\[
  N\otimes_{\mathcal O_S}\Sym(E)(n),
\]
其中 \(N\) 是 \(S\) 上的有限型局部自由层。应用 \(\Proj_0\)，就得到所需的由 \(\xi\) 的幂生成的结论。

其次证明，\(1,\xi,\ldots,\xi^{r+1}\) 在 \(K^\bullet(S)\) 上线性相关。这由典范满同态 \(f^*E\to\mathcal O_X(1)\) 得出。

最后证明，\(1,\xi,\ldots,\xi^r\) 在 \(K^\bullet(S)\) 上线性无关。

\sgaStableTarget{sga6.e06.sec.1.2}\subsection*{1.2. 分次模的定义与性质}
为方便起见，这里回顾分次模的几个经典定义和性质。

设 \(S\) 是概形，\(B\) 是分次 \(\mathcal O_S\)-代数，\(M\) 是分次 \(B\)-模。对每个分次 \(B\)-模 \(N\)，仍用 \(N(n)\) 表示由
\[
  N(n)_k=N_{n+k}
\]
定义的分次模。称 \(M\) 为局部自由分次 \(B\)-模（或有限型局部自由分次 \(B\)-模），是指 \(S\) 的每一点都有一个开邻域 \(U\)，使 \(M|_U\) 作为分次 \(B\)-模同构于若干形如 \(B(n)|_U\) 的模的直和（或有限直和）。称 \(M\) 为有限表现分次 \(B\)-模，是指每一点都有开邻域 \(U\)，使存在正合列 \(P\to Q\to M|_U\to0\)，其中 \(P,Q\) 是有限型自由分次 \((B|_U)\)-模，且同态均与分次相容。

\textbf{1.2.1.} 对每个整数 \(n\)，有分次 \(B\)-模的典范同态
\sgaStableTarget{sga6.e06.eq.1-2-1}\begin{equation}\tag{1.2.1}
  M_n\otimes_{\mathcal O_S}B(-n)\longrightarrow M,
\end{equation}
其定义是 \(x\otimes a\mapsto ax\)。\textbf{1.2.2.} 若分次 \(B\)-模 \(M\) 作为分次模同构于某个 \((B(-n))^k\)，就称 \(M\) 是次数为 \(n\) 的齐次有限型自由模。

设 \(M\) 是有限型自由分次 \(B\)-模。由 1.2，可找到分次同构
\sgaStableTarget{sga6.e06.eq.1-2-2}\begin{equation}\tag{1.2.2}
  M\simeq \bigoplus_{n=n_0}^{n_1} B(-n)^{k_n}.
\end{equation}

这给出 \(M\) 的一个齐次有限型自由分次 \(B\)-模的直和分解
\[
 M=\bigoplus_{n=n_0}^{n_1}M^{(n)},
\]
其中 \(M^{(n)}\) 的次数为 \(n\)，秩为 \(k_n\)。

\sgaStableTarget{sga6.e06.prop.1-2}\textbf{命题 1.2.3.} 分解 \hyperlink{sga6.e06.eq.1-2-2}{(1.2.2)} 中出现的整数 \(k_n\) 与所选分解无关。

设给定两个这样的分解，因而有一个分次同构
\[
 \bigoplus_{n=n_0}^{n_1}(B(-n))^{k_n}
 \simeq \bigoplus_{n=n'_0}^{n'_1}(B(-n))^{k'_n}.
\]
对左边非零齐次自由直和因子的数目作归纳。若没有这样的因子，则 \(M=0\)，命题显然。假设已知因子数为 \(m\) 时成立，来证明因子数为 \(m+1\) 的情形。给定同构与分次相容，因而在两边的齐次分量上给出同构。先看次数 \(<n_0\) 的齐次分量，得到此时 \(k'_n=0\)；再看次数为 \(n_0\) 的齐次分量，它们的秩相同，因此 \(k_{n_0}=k'_{n_0}\)。给定同构由此诱导两分解中最低次齐次自由直和因子的同构，并可除以这些因子而过渡到商。在此商模上应用归纳假设，即得结论。

特别地，由 1.2.3，对于每个有限型自由分次模 \(M\)，齐次基中具有给定次数的元素个数，与基的选择无关。因而，由 \(k_{n_0}\ne0\)、\(k_{n_1}\ne0\) 确定的分解 (1.2.2) 中的两个整数 \(n_0,n_1\)，也与分解无关。

\sgaStableTarget{sga6.e06.prop.1-3}\textbf{命题 1.3.} 设 \(S\) 拟紧，\(E\) 是有限型局部自由 \(\mathcal O_S\)-模，\(X=\mathbf P(E)\)，\(F\) 是一个在 \(S\) 上平坦的有限表现 \(\mathcal O_X\)-模。则存在整数 \(n_0\)，使得对所有 \(n\geq n_0\)，有
\[
\begin{array}{ll}
\text{\rm (i)}& R^i f_*(F(n))=0\quad (i\geq 1),\\[2mm]
\text{\rm (ii)}& f_*(F(n))\text{ 是 }S\text{ 上的有限型局部自由层},\\[2mm]
\text{\rm (iii)}& \displaystyle M=\bigoplus_{n\geq n_0} f_*(F(n))
\text{ 是有限表现分次 }\Sym(E)\text{-模。}
\end{array}
\]

所求各性质在 \(S\) 上是局部的；由于 \(S\) 拟紧，只需取局部所得整数的最大值，便可归结到 \(S\) 仿射的情形。

先假设 \(S\) 诺特。性质 (i) 来自 Serre 定理。取 \(n_0\) 使性质 (i) 成立后，由 EGA III 7.9.10 得到性质 (ii)。由 EGA II 3.4.5，存在有限型拟凝聚分次 \(\Sym(E)\)-模 \(M'\)，使 \(F=\Proj_0(M')\)。另一方面，由 EGA III 2.3.1，\(\Gamma_*(\Proj_0(M'))=\bigoplus_n f_*(F(n))\) 与 \(M'\) 在充分高次上同构。因此存在整数 \(n_0\)，使 \(\bigoplus_{n\ge n_0}f_*(F(n))\) 是有限型 \(\Sym(E)\)-模，利用诺特假设，它也是有限表现的。事实上，对任意 \(n_0\) 此结论都成立，因为由 Serre 定理，各 \(f_*(F(n))\) 都凝聚。于是命题在 \(S\) 诺特时得证。

在这一归结中，使用笛卡尔方块
\[
\begin{tikzcd}[column sep=large,row sep=large]
X \arrow[r,"g"] \arrow[d,"f"'] & X_0 \arrow[d,"f_0"]\\
S \arrow[r,"g_0"'] & S_0 .
\end{tikzcd}
\]
这里 \(S_0\) 是诺特概形，\(X_0=\mathbf P(E_0)\)，且 \(F=g^*(F_0)\)。前面的论证对 \(F_0\) 证明了命题。由 EGA III 7.7.5 和 7.8.5，在上述假设下基变换定理成立，立即推出命题对 \(F\) 成立。

因此，只需证明一般的仿射情形可以由诺特情形经基变换得到。由于 \(S\) 仿射，它是诺特仿射概形的滤过射影极限：若 \(S=\Spec(A)\)，把 \(A\) 写成其有限型 \(\mathbb Z\)-子代数的归纳极限即可。写成 \(S=\varprojlim S_\lambda\)。由 EGA IV 8.5.2 和 8.5.5，存在指标 \(\lambda\) 和有限型局部自由 \(\mathcal O_{S_\lambda}\)-模 \(E_\lambda\)，使 \(E=u_\lambda^*(E_\lambda)\)，其中 \(u_\lambda:S\to S_\lambda\) 为投影。对 \(\mu\ge\lambda\)，令 \(E_\mu=u_{\lambda\mu}^*E_\lambda\)，\(X_\mu=\mathbf P(E_\mu)\)，则 \(X_\mu=X_\lambda\times_{S_\lambda}S_\mu\)，且 \(X=\varprojlim X_\mu\)。所以存在指标 \(\lambda'\ge\lambda\) 及 \(X_{\lambda'}\) 上的有限表现层 \(F_{\lambda'}\)，使其拉回为 \(F\)。由于 \(F_{\lambda'}\) 和相应的射影丛均为有限表现，由 EGA IV 11.2.6，还可增大指标，使 \(F_{\lambda''}\) 在 \(S_{\lambda''}\) 上平坦。此时 1.3 的全部假设都得到满足，证明完成。

\sgaStableTarget{sga6.e06.cor.1-4}\textbf{推论 1.4.} 在\hyperlink{sga6.e06.prop.1-3}{命题 1.3} 的假设下，存在一个有限表现、且在 \(\mathcal O_S\) 上平坦的分次 \(\Sym(E)\)-模 \(M\)，使得
\[
  F=\Proj_0(M).
\]

取一个使\hyperlink{sga6.e06.prop.1-3}{命题 1.3} 的各项性质成立的整数 \(n_0\)。对于
\[
 M=\bigoplus_{n\ge n_0}f_*(F(n)),
\]
我们已经看到，\(M\) 是有限表现的；它在 \(S\) 上平坦，因为当 \(n\ge n_0\) 时，\(f_*(F(n))\) 在 \(S\) 上局部自由。此外，由 EGA II 3.4，\(F=\Proj_0(M)\)。现在转向证明的 b) 部分。

\sgaStableTarget{sga6.e06.lem.1-5-1}\textbf{引理 1.5.1.} 设 \(A\) 是环，\(B\) 是满足 \(B_0=A\) 的正分次 \(A\)-代数，\(M\) 是分次 \(B\)-模。假设把 \(M\) 看作不分次的 \(B\)-模时，它是有限型投射模。设 \(\widetilde B\) 是 \(B\) 在 \(\Spec(A)\) 上定义的代数层，\(\widetilde M\) 是相应的分次 \(\widetilde B\)-模。则在 \(\Spec(A)\) 上局部，\(\widetilde M\) 作为分次模 \hyperlink{sga6.e06.prop.1-2}{(1.2)}，在 \(\widetilde B\) 上是有限型自由模。

事实上，由于 \(M\) 是有限型的，它的次数有下界。此外，\(M\otimes_B A\) 是有限型投射 \(A\)-模，所以只有有限多个非零齐次分量，每个分量都是有限型投射 \(A\)-模。因此，在 \(\Spec(A)\) 上局部化后，可以假设 \(M\otimes_B A\) 是有限型自由 \(A\)-模。最后，\(M\) 的投射性给出
\[
 \Tor_1^B(M,A)=0.
\]
这些性质蕴涵 \(M\) 作为分次模在 \(B\) 上是有限型自由模（Bourbaki, \emph{Algèbre}, ch.~II, \S~11, prop.~7），因而 \(\widetilde M\) 在 \(\widetilde B\) 上是有限型自由模，如所欲证。

\sgaStableTarget{sga6.e06.lem.1-5-2}\textbf{引理 1.5.2.} 设 \(A\) 是环，\(B\) 是有限表现 \(A\)-代数，\(M\) 是有限表现 \(B\)-模。假设 \(B\) 和 \(M\) 在 \(A\) 上平坦，而且对 \(A\) 的每个剩余域 \(k\)，\(B\otimes_A k\)-模 \(M\otimes_A k\) 的投射维数至多为 \(n\)（若 \(B=A[t_0,\ldots,t_r]\)，即属这种情形）。则 \(M\) 有一个长为 \(n\) 的、由有限型投射 \(B\)-模组成的消解：
\[
  0\longrightarrow L_n\longrightarrow\cdots\longrightarrow L_0\longrightarrow M\longrightarrow 0.
\]
这就是 EGA IV, 12.3.2 的局部判别准则。

\sgaStableTarget{sga6.e06.prop.1-6}\textbf{命题 1.6.} 设 \(S\) 拟紧，\(E\) 局部自由且秩为 \(r+1\)，\(B=\Sym(E)\) 带通常的分次，\(M\) 是在 \(S\) 上平坦的有限表现分次 \(B\)-模。则 \(M\) 有一个分次消解
\[
  0\longrightarrow L_{r+1}\longrightarrow\cdots\longrightarrow L_0\longrightarrow M\longrightarrow 0,
\]
其中 \(L_i\) 是有限型局部自由分次 \(B\)-模。

先证明：对每个满足命题假设的分次 \(B\)-模 \(M\)，都存在一个分次满同态 \(L\to M\)，其中 \(L\) 是有限型局部自由分次 \(B\)-模。由于 \(M\) 有限型，存在整数 \(N\)，使 \(M\) 由次数 \(\leq N\) 的齐次分量在 \(B\) 上生成。令
\[
 L=\bigoplus_{i\leq N}M_i\otimes_{\mathcal O_S}B(-i).
\]
利用 (1.2.1) 定义 \(L\to M\)，此同态显然满。

由 1.5.2，\(M\) 在每个仿射开集上都有一个由有限型局部自由 \(B\)-模组成的不分次消解。因此，若已给出一个全局部分消解
\[
 L_i\longrightarrow L_{i-1}\longrightarrow\cdots
 \longrightarrow L_0\longrightarrow M\longrightarrow0,
\]
其中各 \(L_j\) 都是有限型局部自由 \(B\)-模，则 \(L_i\to L_{i-1}\) 的核 \(Z_i\) 仍为有限型 \(B\)-模（I 2.11）。对 \(i+1\) 作同样论证，可知它甚至为有限表现。最后，它在 \(S\) 上平坦。因此它满足与 \(M\) 相同的假设，利用开头的论证，便可逐步构造 \(M\) 的一个全局分次消解，各项均为有限型局部自由分次 \(B\)-模。对 \(i\geq r+1\)，由 1.5.2，\(Z_i\) 在每个仿射开集上都是投射 \(B\)-模；由 1.5.1，它作为分次 \(B\)-模也局部自由。于是，在已构造的分次消解中，用 \(Z_{r+1}\) 代替 \(L_{r+1}\)，把后续各项换成零，就得到所述消解。% 原书第372页此处核指标与截断指标的配合须另作数学复核；正文忠实保留。


最后证明证明安排中的 c)。

\sgaStableTarget{sga6.e06.prop.1-7}\textbf{命题 1.7.} 设 \(S\) 拟紧，\(E\) 有限型局部自由，\(B=\Sym(E)\)，\(L\) 是有限型局部自由分次 \(B\)-模。则 \(L\) 有一个由分次子模 \(L^{(m)}\) 组成的有限递减滤过，使得
\[
\begin{array}{ll}
\text{\rm (i)}& L^{(m)}=L\quad\text{当 }m\leq m_0,\\
\text{\rm (ii)}& L^{(m)}=0\quad\text{当 }m\geq m_1,\\
\text{\rm (iii)}& \operatorname{gr}^m L
  \simeq N_m\otimes_{\mathcal O_S}B(m),
\end{array}
\]
其中 \(N_m\) 在 \(S\) 上有限型局部自由。

由于 \(S\) 拟紧，可以找到整数 \(n_0,n_1\)，使 \(S\) 的每一点附近，\(L\) 的齐次基元素的次数都介于 \(n_0\) 与 \(n_1\) 之间。对 \(N=n_1-n_0\) 作归纳。

若 \(N=0\)，令 \(n=n_0=n_1\)。典范同态 (1.2.1)
\[
 L_n\otimes_{\mathcal O_S}B(-n)\longrightarrow L
\]
是同构。事实上，这可局部检验；由 1.2.3，可以假设 \(L\simeq(B(-n))^k\)，此时结论显然。

若 \(N>0\)，仍考虑典范同态
\[
 L_{n_0}\otimes_{\mathcal O_S}B(-n_0)\longrightarrow L.
\]
局部检验可知，此同态为单同态，其余核 \(L'\) 有限型局部自由，且局部生成元的次数介于 \(n_0+1\) 与 \(n_1\) 之间。事实上，在 \(S\) 的每一点附近，\(L_{n_0}\) 或者为零，或者是 \(L\) 的最低次非零齐次分量。取 (1.2.2) 型直和分解，后一情形下，由 1.2.3，\(L_{n_0}\otimes B(-n_0)\) 同构于 \(L\) 的直和因子 \((B(-n_0))^{k_{n_0}}\)，故得到所述性质。对 \(L'\) 应用归纳假设，立即得到所需结论。

\sgaStableTarget{sga6.e06.sec.1.8}\subsection*{1.8. 由 \(\xi\) 的幂生成}
回到定理 1.1 的证明，来证明 \(K^\bullet(X)\) 作为 \(K^\bullet(S)\)-模由 \(\xi\) 的幂生成。

若 \(F\) 是有限型局部自由 \(\mathcal O_X\)-模，由 1.4，存在有限表现、且在 \(S\) 上平坦的分次 \(B\)-模 \(M\)，使 \(F=\Proj_0(M)\)。由 1.6，\(M\) 有一个由有限型局部自由分次 \(B\)-模组成的分次消解。由此得到正合列
\[
0\longrightarrow\Proj_0(L_{r+1})\longrightarrow\cdots
\longrightarrow\Proj_0(L_0)\longrightarrow\Proj_0(M)\longrightarrow0
\]
以及 \(K^\bullet(X)\) 中的等式
\[
 \cl(F)=\sum_i(-1)^i\cl(\Proj_0(L_i)).
\]
利用 1.7 在各 \(L_i\) 上定义的滤过，得到
\[
 \cl(\Proj_0(L_i))
 =\sum_j\cl\bigl(\Proj_0(N_{ij}\otimes_{\mathcal O_S}B(j))\bigr).
\]
而 \(\Proj_0(N_{ij}\otimes_{\mathcal O_S}B(j))=f^*(N_{ij})(j)\)，所以
\[
 \cl(\Proj_0(L_i))=\sum_j\cl(f^*N_{ij})\,\xi^j.
\]
最终把 \(\cl(F)\) 表示成了 \(\xi\) 的幂的 \(K^\bullet(S)\)-线性组合。

\sgaStableTarget{sga6.e06.sec.1.9}\subsection*{1.9. 关系式}
为了证明线性相关，使用典范满同态 \(f^*E\to\mathcal O_X(1)\)。其核局部自由且秩为 \(r\)。因此
\[
  \lambda^{r+1}(f^*E-\xi)=0.
\]
这里 \(E=\cl(E)\)，\(K^\bullet(X)\) 上的各 \(\lambda^i\) 运算在 V 2.2 中定义。因此
\[
 \sum_{i=0}^{r+1}\lambda^i(f^*E)\lambda^{r+1-i}(-\xi)=0.
\]
由于 \(\mathcal O_X(1)\) 可逆，\(\lambda_t(\xi)=1+\xi t\)，因而 \(\lambda_t(-\xi)=1/(1+\xi t)\)。最终得到关系式
\sgaStableTarget{sga6.e06.eq.1-9-1}\begin{equation}\tag{1.9.1}
  \sum_{i=0}^{r+1}(-1)^i\,\lambda^i(E)\,\xi^{r+1-i}=0,
\end{equation}
最后证明 \(1,\xi,\ldots,\xi^r\) 在 \(K^\bullet(S)\) 上线性无关。\begingroup\renewcommand{\thefootnote}{*}\footnote[1]{此证明归功于 J.-P. Jouanolou。}\endgroup
\sgaStableTarget{sga6.e06.sec.1.10}\subsection*{1.10--1.11. 线性无关性}

\sgaStableTarget{sga6.e06.lem.1-10}\textbf{引理 1.10.} 在\hyperlink{sga6.e06.thm.1-1}{定理 1.1} 的假设下，设 \(F\) 和 \(G\) 是 \(X\) 上的有限型局部自由层，假设它们在 \(K^\bullet(X)\) 中有相同的类。则存在整数 \(n_0\)，使得对所有 \(n\geq n_0\)，有
\[
\begin{array}{ll}
\text{\rm (i)}& f_*(F(n))\text{ 和 }f_*(G(n))\text{ 是 }S\text{ 上的有限型局部自由层},\\
\text{\rm (ii)}& \cl f_*(F(n))=\cl f_*(G(n))\text{ 在 }K^\bullet(S)\text{ 中成立。}
\end{array}
\]

设 \(A\) 是以 \(X\) 上有限型局部自由层的同构类为基的自由阿贝尔群。根据 \(K^\bullet\) 的构造，\(F,G\) 在 \(K^\bullet(X)\) 中具有相同的类，意味着存在有限型局部自由 \(\mathcal O_X\)-模的正合列
\[
 0\longrightarrow L_i\longrightarrow M_i\longrightarrow N_i
 \longrightarrow0\qquad(i=1,\ldots,k)
\]
以及整数 \(n_i\)，使在 \(A\) 中有
\[
 (F)-(G)=\sum_i n_i\bigl((L_i)-(M_i)+(N_i)\bigr).
\]
对每个 \(n\)，与 \(\mathcal O_X(n)\) 张量定义 \(A\) 的自同构，因此
\[
 (F(n))-(G(n))
 =\sum_i n_i\bigl((L_i(n))-(M_i(n))+(N_i(n))\bigr),
\]
且各列 \(0\to L_i(n)\to M_i(n)\to N_i(n)\to0\) 仍正合。

设 \(B'\) 是以 \(S\) 上拟凝聚层的同构类为基的自由阿贝尔群，\(A'\subset B'\) 是由 \(S\) 上有限型局部自由层的类生成的子群。取正像定义同态 \(A\to B'\)，所以仍有
\[
\begin{aligned}
 &(f_*F(n))-(f_*G(n))\\
 &\quad=\sum_i n_i\bigl((f_*L_i(n))-(f_*M_i(n))+(f_*N_i(n))\bigr).
\end{aligned}
\]
当 \(n\) 充分大时，各 \(R^k f_*L_i(n)\)（\(k\geq1\)）消失，且 \(f_*L_i(n)\)、\(f_*M_i(n)\)、\(f_*N_i(n)\)、\(f_*F(n)\)、\(f_*G(n)\) 均有限型局部自由。于是这些元素全部属于 \(A'\)，满足上面的关系，而且各列
\[
 0\longrightarrow f_*L_i(n)\longrightarrow f_*M_i(n)
 \longrightarrow f_*N_i(n)\longrightarrow0
\]
正合。因此 \(f_*F(n)\) 和 \(f_*G(n)\) 在 \(K^\bullet(S)\) 中具有相同的类。

假设
\sgaStableTarget{sga6.e06.eq.1-11-1}\begin{equation}\tag{1.11.1}
  f^*(x_0)+f^*(x_1)\xi+\cdots+f^*(x_r)\xi^r=0 .
\end{equation}
把每个 \(x_j\) 写成两个层的类之差：\(x_j=\cl(E_j)-\cl(E'_j)\)，其中 \(E_j,E'_j\) 均为 \(S\) 上的有限型局部自由层。令
\[
 F=\bigoplus_{j=0}^r f^*E_j(j),\qquad
 F'=\bigoplus_{j=0}^r f^*E'_j(j).
\]
等式 (1.11.1) 就成为 \(\cl(F)=\cl(F')\)。由引理 1.10，对充分大的 \(n\)，有 \(\cl(f_*F(n))=\cl(f_*F'(n))\)。而
\[
\begin{aligned}
 f_*F(n)&=\bigoplus_j f_*\bigl(f^*E_j\otimes\mathcal O_X(n+j)\bigr)\\
        &=\bigoplus_j E_j\otimes f_*\mathcal O_X(n+j).
\end{aligned}
\]
又 \(f_*\mathcal O_X(n+j)=\Sym^{n+j}(E)\)；用 \(\sigma_{n+j}(E)\) 表示其在 \(K^\bullet(S)\) 中的类。于是 (1.11.1) 最终给出
\sgaStableTarget{sga6.e06.eq.1-11-2}\begin{equation}\tag{1.11.2}
  \sum_{i=0}^{r} x_i\,\sigma_{n+i}(E)=0
  \quad\text{在 }K^\bullet(S)\text{ 中},
\end{equation}
其中 \(\sigma_p(E)\) 表示 \(\Sym^p(E)\) 的类，并约定当 \(p<0\) 时，\(\sigma_p(E)=0\)。

考虑复形
\[
\begin{aligned}
 0\longrightarrow\bigwedge^{r+1}E\otimes\Sym(E)
 \longrightarrow\cdots\longrightarrow E\otimes\Sym(E)
 \longrightarrow\Sym(E)\longrightarrow\mathcal O_S\longrightarrow0,
\end{aligned}
\]
其中箭头由
\[
\begin{aligned}
 &(x_1\wedge\cdots\wedge x_n)\otimes a\\
 &\quad\longmapsto\sum_{j=1}^n(-1)^{j-1}
 (x_1\wedge\cdots\wedge\widehat{x_j}\wedge\cdots\wedge x_n)\otimes ax_j
\end{aligned}
\]
定义。把 \(\bigwedge^n E\) 看作集中于次数 \(n\) 的分次 \(\mathcal O_S\)-模，这就是分次 \(\Sym(E)\)-模的复形，其态射均为分次态射。

这个复形是 \(\mathcal O_S\) 的消解：在 \(E\) 自由的开集上，\(\Sym(E)\) 等同于多项式代数 \(A[T_0,\ldots,T_r]\)，此复形等同于由 \(\Sym(E)\)-正则序列 \(T_0,\ldots,T_r\) 定义的 Koszul 复形；它是 \(\Sym(E)/(T_0,\ldots,T_r)=\mathcal O_S\) 的消解。

写出各正次数齐次分量无环，便在 \(K^\bullet(S)\) 中得到：对所有 \(p\geq1\)，
\sgaStableTarget{sga6.e06.eq.1-11-3}\begin{equation}\tag{1.11.3}
  \sum_{i=0}^{r+1}(-1)^i\lambda^i(E)\,\sigma_{p-i}(E)=0 .
\end{equation}
假设 (1.11.2) 对 \(n\geq n_0\) 成立。来证明它对 \(n_0-1\) 仍成立；由此可以递降归纳到 \(n=-r\)，后面将看到这足以推出各 \(x_i\) 为零。

由 (1.11.3)，对任意 \(p(j)\geq1\)，有
\[
 \sum_{j=0}^r x_j\left(\sum_{i=0}^{r+1}
 (-1)^i\lambda^i(E)\sigma_{p(j)-i}(E)\right)=0.
\]
取 \(p(j)=n_0+r+j\)；当 \(n_0>-r\) 时，它确实 \(\geq1\)。得到
\sgaStableTarget{sga6.e06.eq.1-11-4}\begin{equation}\tag{1.11.4}
\begin{aligned}
 &\sum_{j=0}^r x_j\sum_{i=0}^{r+1}(-1)^i\lambda^i(E)
                  \sigma_{n_0+r+j-i}(E)\\
 &=\sum_{i=0}^{r+1}(-1)^i\lambda^i(E)
                  \sum_{j=0}^r x_j\sigma_{n_0+r+j-i}(E)=0.
\end{aligned}
\end{equation}
若 \(i\ne r+1\)，则 \(n_0+r-i\geq n_0\)，由 (1.11.2)，相应的内层和为零。因此 (1.11.4) 化为
\[
 \lambda^{r+1}(E)\sum_{j=0}^r x_j\sigma_{n_0-1+j}(E)=0.
\]
\(E\) 的秩为 \(r+1\)，故 \(\lambda^{r+1}(E)\) 是可逆层的类，在 \(K^\bullet(S)\) 中可逆。因此得到 \(n=n_0-1\) 时的 (1.11.2)。

所以 (1.11.2) 对 \(n\geq-r\) 成立。取 \(n=-r\)，得到 \(x_r\sigma_0(E)=0\)。由于 \(\sigma_0(E)=1\)，有 \(x_r=0\)。再使 \(n\) 从 \(-r\) 依次取到 \(0\)，对 \(j\) 作递降归纳，就得到所有 \(x_j=0\)。这证明了 \(1,\xi,\ldots,\xi^r\) 在 \(K^\bullet(S)\) 上线性无关，完成了定理 1.1 的证明。

\sgaStableTarget{sga6.e06.rem.1-12}\textbf{注 1.12.} 对每个有限型局部自由 \(\mathcal O_S\)-模 \(E\)，设 \(\sigma_i(E)\) 是 \(E\) 的第 \(i\) 个对称幂在 \(K^\bullet(S)\) 中的类。令
\[
  \sigma_t(E)=\sum_{i\geq 0}\sigma_i(E)t^i .
\]
关系式 \hyperlink{sga6.e06.eq.1-11-3}{(1.11.3)} 等价于
\[
  \lambda_t(E)\,\sigma_{-t}(E)=1.\tag{1.12.1}
\]
特别地，\(\sigma_t\) 是可加的，并延拓为同态
\[
  K^\bullet(S)\longrightarrow 1+K^\bullet(S)[[t]]_+ .
\]

\sgaStableTarget{sga6.e06.rem.1-13}\textbf{注 1.13.} 可以取 \(\xi^k,\xi^{k+1},\ldots,\xi^{k+r+1}\) 作为 \(K^\bullet(X)\) 在 \(K^\bullet(S)\) 上的一组基。\footnote{译校注：法文原书第 379 页的末项印作 \(\xi^{k+r+1}\)；与定理 1.1 所述秩及本段论证相符的末项应为 \(\xi^{k+r}\)，即连续 \(r+1\) 个幂。}因为 \(\xi\) 可逆，这些元素线性无关；另一方面，由 (1.9.1) 及 \(\xi\) 的可逆性，其他各幂都可以显式表示成它们的 \(K^\bullet(S)\)-线性组合。

令 \(x=1-\xi^{-1}\)（\(x\) 可解释为“超平面截面的类”）。于是也可取 \(1,x,\ldots,x^r\) 为基。为得到 \(1,x,\ldots,x^{r+1}\) 的线性相关方程，考虑正合列
\[
 0\longrightarrow F\longrightarrow f^*E\longrightarrow\mathcal O_X(1)\longrightarrow0,
\]
以及取对偶所得的正合列
\[
 0\longrightarrow\mathcal O_X(-1)\longrightarrow f^*\check E
 \longrightarrow\check F\longrightarrow0.
\]
注意 \(\gamma^{r+1}(\check F-r)=0\)，从而
\[
 \gamma^{r+1}(\check F-r)
 =\gamma^{r+1}(\check E-r-1+x)=0,
\]
即
\sgaStableTarget{sga6.e06.eq.1-13-1}\begin{equation}\tag{1.13.1}
 \sum_{k=0}^{r+1}\gamma^{r+1-k}(\check E-r-1)x^k=0.
\end{equation}
引入 \(x'=\xi-1\) 也得到类似结果。

最后注意，若 \(E=\mathcal O_S^{r+1}\)，从而 \(X=\mathbf P_S^r\)，则 \(E=r+1\)，等式 (1.13.1) 化为
\sgaStableTarget{sga6.e06.eq.1-13-2}\begin{equation}\tag{1.13.2}
 x^{r+1}=0.
\end{equation}
\sgaStableTarget{sga6.e06.rem.1-14}\textbf{注 1.14.} 使用 Hakim 意义下的相对概形，当 \(S\) 是局部赋环拓扑斯时，\hyperlink{sga6.e06.thm.1-1}{定理 1.1} 仍然有意义；若该拓扑斯的终对象拟紧，证明也可以直接移用而无实质变化。

\sgaStableTarget{sga6.e06.sec.2}\section*{2. 射影丛的 \(K^\bullet\) 的计算：完美复形的情形}

本节不用于讨论班的其余部分。本节中，记号 \(K^\bullet(X)\) 表示 \(X\) 上完美复形的 Grothendieck 群（\hyperlink{sga6.e04.def.2-2}{IV 2.2}）。我们将证明下述与 \hyperlink{sga6.e06.thm.1-1}{1.1} 类似的定理：

\sgaStableTarget{sga6.e06.thm.2-1}\textbf{定理 2.1.} 设 \(S\) 拟紧，\(E\) 是秩为 \(r+1\) 的局部自由 \(\mathcal O_S\)-模，\(X=\mathbf P(E)\) 是相应的射影丛；设 \(\xi\) 是在零次等于 \(\mathcal O_X(1)\)、其余次数为零的复形在 \(K^\bullet(X)\) 中的类。则带有自然 \(K^\bullet(S)\)-代数结构的 \(K^\bullet(X)\)，是以
\[
  1,\xi,\ldots,\xi^r .
\]
为基的自由 \(K^\bullet(S)\)-模。
\sgaStableTarget{sga6.e06.sec.2.2}\subsection*{2.2. 记号}
固定如下记号：\(B\) 表示 \(E\) 的对称代数；\(\mathcal A\) 表示 \(\mathcal O_X\)-模范畴，\(\mathcal G\) 表示以分次态射为态射的分次 \(B\)-模范畴，\(\mathcal G_0\) 表示 \(\mathcal G\) 的如下满子范畴：其对象是作为分次 \(B\)-模有限型局部自由的 \(B\)-模；\(\mathcal G'\) 表示 \(\mathcal O_S\)-模范畴。范畴 \(\mathcal A,\mathcal G,\mathcal G'\) 都是阿贝尔范畴。对每个 \(n\) 以及 \(D^+(\mathcal A)\) 的每个复形 \(L^\bullet\)，我们用 \(L^\bullet(n)\) 表示复形
\[
  L^\bullet(n)=L^\bullet\otimes_{\mathcal O_X}\mathcal O_X(n).
\]

\sgaStableTarget{sga6.e06.sec.2.3}\subsection*{2.3. 函子 \(\Proj_0\)}
我们需要把函子 \(\Proj_0\) 应用于任意分次 \(B\)-模，而不只限于 EGA II 3.2 的构造中的拟凝聚分次 \(B\)-模。因此，我们简要说明如何在一般情形下作此构造。

\sgaStableTarget{sga6.e06.para.2-3-1}\textbf{2.3.1.} 首先，设 \(S\) 是概形，\(A\) 是拟凝聚 \(\mathcal O_S\)-代数，\(Y=\Spec(A)\)，\(f:Y\to S\) 是结构态射。由同构
\[
  A\simeq f_*\mathcal O_Y
\]
得到典范同态
\sgaStableTarget{sga6.e06.eq.2-3-1-1}\begin{equation}\tag{2.3.1.1}
  f^{-1}(A)\longrightarrow \mathcal O_Y,
\end{equation}
其中 \(f^{-1}\) 表示集合层意义下的逆像；换言之，有赋环空间态射
\sgaStableTarget{sga6.e06.eq.2-3-1-2}\begin{equation}\tag{2.3.1.2}
  \bar f:(Y,\mathcal O_Y)\longrightarrow (S,A).
\end{equation}
它给出一个从 \(A\)-模范畴到 \(\mathcal O_Y\)-模范畴的逆像函子 \(\bar f^{\,*}\)，这是从 \(\mathcal O_Y\)-模范畴到 \(f_*(\mathcal O_Y)\)-模范畴的正像函子 \(\bar f_*\) 的左伴随。因此，对每个 \(A\)-模 \(M\)，有
\sgaStableTarget{sga6.e06.eq.2-3-1-3}\begin{equation}\tag{2.3.1.3}
  \bar f^{\,*}(M)=f^{-1}(M)\otimes_{f^{-1}A}\mathcal O_Y
\end{equation}


由于同态 \hyperlink{sga6.e06.eq.2-3-1-1}{(2.3.1.1)} 平坦（在每个茎上它对应于局部化），刚才定义的函子 \(\bar f^{\,*}\) 正合（参见 [4]，第 164 页）。若 \(M\) 是拟凝聚 \(A\)-模，此构造给出的就是 EGA II 1.4.3 中定义的模 \(\widetilde M\)。事实上，由同构
\[
  M\simeq f_*(\widetilde M)=\bar f_*(\widetilde M)
\]
通过伴随得到同态
\[
  \bar f^{\,*}(M)\longrightarrow \widetilde M
\]
并且逐茎立即可验证它是同构。

\sgaStableTarget{sga6.e06.para.2-3-2}\textbf{2.3.2.} 现在设 \(B\) 是拟凝聚分次 \(\mathcal O_S\)-代数，\(M\) 是分次 \(B\)-模，\(X=\Proj(B)\)。对 \(S\) 的每个仿射开集 \(U\)，以及 \(B\) 在 \(U\) 上每个次数为 \(d>0\) 的齐次截面 \(s\)，用 \(M_{(s),U}\) 表示分式模 \(M_s\) 的 \(0\) 次齐次分量；于是
\[
  M_{(s),U}=
  \bigl((M|_U)\otimes_{B|_U}(B|_U)_s\bigr)_0 .
\]
若 \(V\) 是 \(U\) 的仿射开集，则 \(M_{(s),V}\) 是 \(M_{(s),U}\) 在 \(V\) 上的限制。

考虑由如下仿射开集形成的 \(X\) 的开集基：
\[
  D_+(U,s)\simeq \Spec\bigl((B|_U)_{(s)}\bigr)
\]
用 \(f\) 表示结构态射 \(X\to S\)，用 \(f_{U,s}\) 表示仿射态射
\[
  f_{U,s}:D_+(U,s)\longrightarrow U .
\]
令
\[
  \Proj_0(M)_{U,s}
  =
  \bar f_{U,s}^{\,*}\bigl(M_{(s),U}\bigr).
\]

使用 \hyperlink{sga6.e06.para.2-3-1}{2.3.1} 的记号，我们将证明，当 \(U\) 和 \(s\) 变化时，各 \(\Proj_0(M)_{U,s}\) 能够粘合，从而定义一个记为 \(\Proj_0(M)\) 的 \(\mathcal O_X\)-模。固定 \(s\)、使 \(U\) 变化时的粘合是显然的，因而可归结为（参见 EGA II 2.5.2）证明：若 \(t\) 是 \(B\) 在 \(U\) 上的次数为 \(e>0\) 的齐次截面，则 \(\Proj_0(M)_{U,s}\) 在 \(D_+(U,st)\) 上的限制典范同构于 \(\Proj_0(M)_{U,st}\)。考虑
\[
\begin{tikzcd}[column sep=large,row sep=large]
D_+(st) \arrow[r,"i"] \arrow[dr,"f_{U,st}"'] &
D_+(s) \arrow[d,"f_{U,s}"]\\
& U
\end{tikzcd}
\]
我们要证明模
\[
  i^*\bar f_{U,s}^{\,*}\bigl(M_{(s),U}\bigr)
  \quad\text{与}\quad
  \bar f_{U,st}^{\,*}\bigl(M_{(st),U}\bigr)
\]
典范同构。一方面，有典范同构
\[
\begin{aligned}
i^*\bar f_{U,s}^{\,*}\bigl(M_{(s),U}\bigr)
&\simeq
i^*\Bigl(f_{U,s}^{-1}\bigl(M_{(s),U}\bigr)
 \otimes_{f_{U,s}^{-1}((B|_U)_{(s)})}\mathcal O_{D_+(s)}\Bigr)\\
&\simeq
i^{-1}\Bigl(f_{U,s}^{-1}\bigl(M_{(s),U}\bigr)
 \otimes_{f_{U,s}^{-1}((B|_U)_{(s)})}\mathcal O_{D_+(s)}\Bigr)
 \otimes_{i^{-1}(\mathcal O_{D_+(s)})}\mathcal O_{D_+(st)}\\
&\simeq
f_{U,st}^{-1}\bigl(M_{(s),U}\bigr)
 \otimes_{f_{U,st}^{-1}((B|_U)_{(s)})}\mathcal O_{D_+(st)} .
\end{aligned}
\]
另一方面，
\[
  \bar f_{U,st}^{\,*}\bigl(M_{(st),U}\bigr)
  \simeq
  f_{U,st}^{-1}\bigl(M_{(st),U}\bigr)
  \otimes_{f_{U,st}^{-1}((B|_U)_{(st)})}\mathcal O_{D_+(st)} .
\]
最后，把 EGA II 2.2.2 译为层的语言，便得到典范同构
\[
  M_{(st),U}\simeq
  \bigl(M_{(s),U}\bigr)_{t^d/s^e}
  \simeq
  M_{(s),U}\otimes_{(B|_U)_{(s)}}(B|_U)_{(st)},
\]
由此即得所需同构。

显然，若 \(M\) 拟凝聚，则如此构造的模 \(\Proj_0(M)\) 确实就是 EGA II 中定义的模。

\sgaStableTarget{sga6.e06.para.2-3-3}\textbf{2.3.3.} 由 \hyperlink{sga6.e06.para.2-3-1}{2.3.1} 和 \hyperlink{sga6.e06.para.2-3-2}{2.3.2} 可知，函子 \(\Proj_0\) 正合。回到 \hyperlink{sga6.e06.sec.2.2}{2.2} 的记号，过渡到导出范畴，便得到一个仍记为 \(\Proj_0\) 的函子：
\[
  \Proj_0:D^+(\mathcal G)\longrightarrow D^+(\mathcal A).
\]
从上述构造也可清楚看出，分次 \(B\)-模范畴上的函子 \(\Proj_0\) 与张量积相容；在导出范畴上所得的函子仍有此性质。

\sgaStableTarget{sga6.e06.sec.2.4}\subsection*{2.4. 函子 \(F_k\)}
对 \(\mathcal O_X\)-模 \(L\) 和整数 \(k\)，令
\sgaStableTarget{sga6.e06.eq.2-4-1}\begin{equation}\tag{2.4.1}
  F_k(L)=\bigoplus_{n\geq k} f_*(L(n)).
\end{equation}
函子 \(F_k:\mathcal A\to\mathcal G\) 左正合。由于 \(\mathcal O_X\)-模范畴有足够的内射对象，可以定义右导出函子
\sgaStableTarget{sga6.e06.eq.2-4-2}\begin{equation}\tag{2.4.2}
  RF_k:D^+(\mathcal A)\longrightarrow D^+(\mathcal G).
\end{equation}
若 \(L\) 是有限型局部自由 \(\mathcal O_X\)-模，则存在整数 \(k_0\)，使对 \(k\geq k_0\)，\(L\) 为 \(F_k\)-无环。只需证明，对充分大的 \(n\)，\(L\) 关于函子 \(L\mapsto f_*(L(n))\) 无环。由于与 \(\mathcal O_X(n)\) 张量是正合函子，这立即由复合函子谱序列及 1.3 得到。

\subsection*{2.5}
我们要对每个 \(\mathcal O_X\)-模 \(F\)，构造一个关于 \(F\) 具有函子性的同态
\sgaStableTarget{sga6.e06.eq.2-5-1}\begin{equation}\tag{2.5.1}
 \beta:\Proj_0(F_k(F))\longrightarrow F.
\end{equation}
为此，定义 \(\beta\) 在各 \(D_+(U,s)\)（2.3.2）上的限制。由构造，令 \(M=F_k(F)\)，则
\[
 \Proj_0(M)|_{D_+(U,s)}=\bar f_{U,s}^{\,*}(M_{(s),U}).
\]
为了定义同态
\[
 \bar f_{U,s}^{\,*}(M_{(s),U})\longrightarrow F|_{D_+(U,s)},
\]
根据 \(\bar f_{U,s}^{\,*}\) 是正像函子的左伴随这一定义，只需定义一个 \(f_{U,s*}(\mathcal O_X)\)-线性同态
\[
 M_{(s),U}\longrightarrow f_{U,s*}(F|_{D_+(U,s)}).
\]
为此，需对 \(U\) 的每个开集 \(V\) 给出同态
\sgaStableTarget{sga6.e06.eq.2-5-2}\begin{equation}\tag{2.5.2}
 \Gamma(V,M)_{(s)}\longrightarrow\Gamma(D_+(V,s),F),
\end{equation}
使其与 \(U\) 中开集的限制相容，再取相应的伴随层同态。具体定义如下：设 \(z/s^n\in\Gamma(V,M)_{(s)}\)。若 \(s\) 的次数为 \(d\)，则
\[
 z\in\Gamma(V,M_{nd})=\Gamma(f^{-1}(V),F(nd)).
\]
把 \(z/s^n\) 送到 \(F\) 在 \(D_+(V,s)\) 上的截面 \(z\,q_d(s^n)^{-1}\)，这里
\[
 q_e:\Gamma(V,B_e)\longrightarrow\Gamma(f^{-1}(V),\mathcal O_X(e))
\]
是典范同态（参见 EGA II 2.6.3 和 2.6.4）。容易验证所需的相容性，由粘合得到 (2.5.1)。\par\smallskip\noindent\textit{译校注：}法文原书第 385 页印作 \(q_d(s^n)\)；按下行 \(q_e\) 的定义及 \(s^n\) 的次数，此处的下标应为 \(nd\)。

显然，若 \(F\) 拟凝聚，如此构造的同态就是 EGA II (3.3.4.1) 中的同态。

\(\mathcal O_X\)-模范畴上的函子同态 \(\Proj_0\circ F_k\to\operatorname{id}\)，过渡到导出范畴后，给出 \(D^+(\mathcal A)\) 上的函子同态
\sgaStableTarget{sga6.e06.eq.2-5-3}\begin{equation}\tag{2.5.3}
 \Proj_0\circ RF_k\longrightarrow\operatorname{id}_{D^+(\mathcal A)}.
\end{equation}

\sgaStableTarget{sga6.e06.prop.2-6}\textbf{命题 2.6.} 设 \(S\) 拟紧，\(E\) 有限型局部自由，\(X=\mathbf P(E)\)，\(L^\bullet\) 是 \(X\) 上的完美复形。则当 \(k\) 足够大时，
\[
\begin{array}{ll}
\text{\rm (i)}& RF_k(L^\bullet)\text{ 是 }\mathcal G_0\text{-完美的},\\[1mm]
\text{\rm (ii)}& \Proj_0 RF_k(L^\bullet)\longrightarrow L^\bullet
\text{ 是 }D^+(\mathcal A)\text{ 中的同构。}
\end{array}
\]

这里所用的 \(\mathcal G_0\)-完美复形概念是局部概念，即关于 \(\mathcal G_0\) 在 \(S\) 上定义的纤维范畴而言的完美性。此概念确实有意义，因为 \(\mathcal G_0\) 满足 I 4.0 的公理：只有对满同态的核的稳定性并非立即可见，而这一点由 1.5.1 得到。

由于 \(S\) 拟紧，只需在 \(S\) 上局部验证 2.6 的两个断言，所以可假设 \(S\) 仿射。此时可逆层 \(\mathcal O_X(1)\) 丰沛，在 \(X\) 上有全局消解（II 2.2.9）。因此存在严格完美复形 \(L'{}^\bullet\)，以及完美复形的导出范畴中的同构 \(L^\bullet\simeq L'{}^\bullet\)，故可假设 \(L^\bullet\) 严格完美。由 2.4，当 \(k\) 足够大时，各 \(L^i\) 都为 \(F_k\)-无环；除有限多个以外，这些项均为零。因此对充分大的 \(k\)，
\[
 RF_k(L^\bullet)=F_k(L^\bullet).
\]
由 1.3 和 1.6，可找到 \(k_0\)，使对 \(k\geq k_0\)，上一条件成立，而且每个 \(F_k(L^i)\) 都是 \(\mathcal G_0\)-完美 \(B\)-模。于是复形 \(F_k(L^\bullet)\) 为 \(\mathcal G_0\)-完美，得到 i)。

如此选取 \(k\) 后，同态 (2.5.3) 就是
\[
 \Proj_0(F_k(L^\bullet))\longrightarrow L^\bullet.
\]
由于已假设 \(L^\bullet\) 严格完美，它在每个分量上都是同构，得到 ii)。

\sgaStableTarget{sga6.e06.sec.2.7}\subsection*{2.7. \(\mathcal G_0\)-完美复形的分解}
在 \(\mathcal G\) 上，把一个分次 \(B\)-模送到其第 \(n\) 个齐次分量的函子，是取值于 \(\mathcal G'\) 的正合函子；它定义函子
\sgaStableTarget{sga6.e06.eq.2-7-1}\begin{equation}\tag{2.7.1}
  \Comp_n:D^+(\mathcal G)\longrightarrow D^+(\mathcal G').
\end{equation}
对每个整数 \(n\)，还可考虑由与 \(B(-n)\) 张量定义的函子 \(\mathcal G'\to\mathcal G\)。同态 \hyperlink{sga6.e06.eq.1-2-1}{(1.2.1)} 先给出 \(\mathcal G\) 上的函子态射，再给出相应导出范畴间的函子态射
\sgaStableTarget{sga6.e06.eq.2-7-2}\begin{equation}\tag{2.7.2}
  \Comp_n\otimes_{\mathcal O_S}B(-n)\longrightarrow
  \operatorname{id}_{D^+(\mathcal G)}.
\end{equation}

设 \(M^\bullet\) 是 \(\mathcal G_0\)-完美复形。我们将证明，存在有限个可区别三角形组成的序列
\[
\begin{array}{ccccccccccc}
& P_1^\bullet &&&& P_2^\bullet &&&& P_k^\bullet &\\
\swarrow && \nwarrow && \swarrow && \nwarrow && \swarrow && \nwarrow\\
N_1^\bullet & \longrightarrow & M^\bullet, & \qquad &
N_2^\bullet & \longrightarrow & P_1^\bullet, & \cdots, &
N_k^\bullet & \longrightarrow & P_{k-1}^\bullet,
\end{array}
\]
其中最后一步的 \(P_k^\bullet\) 是无环的，各项 \(N_i^\bullet\) 形如
\[
  N_i^\bullet=K_i^\bullet\otimes_{\mathcal O_S}B(d_i),
\]
而 \(K_i^\bullet\) 是 \(S\) 上通常意义下的完美复形。

由于 \(M^\bullet\) 是 \(\mathcal G_0\)-完美的，且 \(S\) 拟紧，存在 \(S\) 的一个有限覆盖及整数 \(d_0\)，使在此覆盖的每个开集 \(U\) 上，\(M^\bullet\) 同构于一个由有限型局部自由 \(B\)-模组成的有界复形 \(M_U^\bullet\)，且当 \(d<d_0\) 时，\(\Comp_d(M_U^\bullet)=0\)。

考虑典范态射 (2.7.2)
\[
 N_1^\bullet=\Comp_{d_0}(M^\bullet)\otimes_{\mathcal O_S}B(-d_0)
 \longrightarrow M^\bullet,
\]
设 \(P_1^\bullet\) 为其锥。它仍为 \(\mathcal G_0\)-完美复形。事实上，对每个 \(n\)，\(\Comp_n(M^\bullet)\) 是 \(S\) 上的完美复形，因为当 \(M^\bullet\) 有界且各项为有限型局部自由 \(B\)-模时，这一点成立。因此 \(\Comp_{d_0}(M^\bullet)\otimes B(-d_0)\) 是 \(\mathcal G_0\)-完美复形，其锥也如此（I 4.10）。

另一方面，在覆盖的每个开集 \(U\) 上可写出交换图
\[
\begin{tikzcd}[column sep=large,row sep=large]
 \Comp_{d_0}(M^\bullet)\otimes B(-d_0) \arrow[r] \arrow[d,"\sim"'] &
 M^\bullet \arrow[d,"\sim"]\\
 \Comp_{d_0}(M_U^\bullet)\otimes B(-d_0) \arrow[r] & M_U^\bullet.
\end{tikzcd}
\]
它诱导两个锥之间的同构。而下方同态在次数 \(d_0\) 上为同构，所以 \(P_1^\bullet|_U\) 同构于一个有界复形，其各项为有限型局部自由 \(B\)-模，且 \(d<d_0+1\) 时的齐次分量 \(\Comp_d\) 均为零。因此 \(P_1^\bullet\) 满足与 \(M^\bullet\) 相同的性质，只是 \(d_0\) 换成了 \(d_0+1\)。

可以迭代此过程，只需验证有限步后得到的复形 \(P_k^\bullet\) 无环。这仍是局部问题，故只需对 \(M_U^\bullet\) 验证。对它的每一项，上述态射恰好是定义命题 1.7 的滤过的那些态射。该滤过有限，所以从 \(M_U^\bullet\) 出发经过有限步就得到零复形，即相应的 \(P_k^\bullet\) 无环。

\sgaStableTarget{sga6.e06.sec.2.8}\subsection*{2.8. 定理 2.1 的证明}
设 \(L^\bullet\) 是 \(X\) 上的完美复形。存在整数 \(k\)，使 \(M^\bullet=RF_k(L^\bullet)\) 为 \(\mathcal G_0\)-完美复形，且 \(L^\bullet\simeq\Proj_0(RF_k(L^\bullet))\)。对 \(M^\bullet\) 应用 2.7，再应用函子 \(\Proj_0\)，所构造的可区别三角形序列给出完美复形的导出范畴中的可区别三角形序列；显然，\(\Proj_0\) 把 \(\mathcal G_0\)-完美复形变为完美复形。此外，对 \(D^+(\mathcal G')\) 中的 \(K^\bullet\)，有
\[
 \Proj_0\bigl(K^\bullet\otimes_{\mathcal O_S}B(d)\bigr)=f^*(K^\bullet)(d).
\]
注意，由于 \(f\) 平坦，\(f^*=Lf^*\)。因此，在 \(K^\bullet(X)\) 中得到
\[
 \cl(L^\bullet)=\sum_i\cl(f^*K_i^\bullet)\xi^{d_i}
              =\sum_i f^*\cl(K_i^\bullet)\xi^{d_i},
\]
其中 \(K_i^\bullet\) 是 2.7 中的 \(S\) 上完美复形。所以，\(K^\bullet(X)\) 作为 \(K^\bullet(S)\)-模由 \(\xi\) 的幂生成。

为证明 \(1,\xi,\ldots,\xi^{r+1}\) 线性相关，考虑同态
\[
 K^\bullet_{\mathrm{naive}}(X)\longrightarrow K^\bullet_{\mathrm{parf}}(X),
\]
它把有限型局部自由层 \(L\) 的类，送到在 \(0\) 次为 \(L\)、其余次数为零的复形的类。方程 (1.9.1) 由此给出 \(K^\bullet_{\mathrm{parf}}(X)\) 中的线性相关关系。

为证明 \(1,\xi,\ldots,\xi^r\) 线性无关，可以使用 \(K^\bullet(S)\)-线性态射
\[
 f_*:K^\bullet_{\mathrm{parf}}(X)\longrightarrow K^\bullet_{\mathrm{parf}}(S).
\]
有 \(f_*(1)=1\)，因为 \(f_*\mathcal O_X=\mathcal O_S\)，且 \(R^if_*\mathcal O_X=0\)（\(i\geq1\)）。还具有 \(f_*(\xi^n)=0\)（\(-r\leq n<0\)），因为此时所有 \(R^if_*\mathcal O_X(n)\)（\(i\geq0\)）都为零。

若有关系式
\[
 f^*(x_0)+f^*(x_1)\xi+\cdots+f^*(x_r)\xi^r=0,
\]
乘以 \(\xi^{-r}\) 并应用 \(f_*\)，得到 \(x_r=0\)。由此归纳即可：若 \(0\leq s\leq r\)，且对 \(s<j\leq r\) 有 \(x_j=0\)，则
\[
 \sum_{i=0}^s f^*(x_i)\xi^i=0,
\]
从而乘以 \(\xi^{r-s}\) 后有
\[
 \sum_{i=0}^s f^*(x_i)\xi^{i+r-s}=0,
\]
由前面的论证得到 \(x_s=0\)。这完成了证明。

% 第 VI 讲的续接宏
\providecommand{\Sym}{\operatorname{Sym}}
\providecommand{\GL}{\operatorname{GL}}
\providecommand{\Grass}{\operatorname{Grass}}
\providecommand{\gr}{\operatorname{gr}}
\providecommand{\ctop}{c_{\mathrm{top}}}
\providecommand{\Rep}{\operatorname{Rep}}
\providecommand{\RZ}{R_{\mathbb Z}}
\providecommand{\Dflag}{D}
\providecommand{\Gmn}{G}


\sgaStableTarget{sga6.e06.sec.3}\section*{3. 射影丛的 \(K^\bullet\) 的结构定理的推论}

\sgaStableTarget{sga6.e06.prop.3-1}\textbf{命题 3.1.} 设 \(S\) 是拟紧概形，\(E_1,\ldots,E_n\) 是有限型局部自由 \(\mathcal O_S\)-模，\(D(E_i)\) 是 \(E_i\) 的旗概形。令
\[
  X=D(E_1)\times_S\cdots\times_S D(E_n).
\]
则典范同态
\[
  K^\bullet(S)\longrightarrow K^\bullet(X)
\]
为单同态，其中 \(K^\bullet\) 可以用有限型局部自由 \(\mathcal O_S\)-模定义，也可以用完美复形定义。

利用旗概形的泛性质，立即可验证这一点。概形 \(X\) 可如下构造：考虑一列归纳定义的 \(S\)-概形 \(X_1,\ldots,X_n\)，取 \(X_k\) 为 \(E_k\) 在 \(X_{k-1}\) 上的逆像的旗概形；于是 \(X=X_n\)。因此，为了证明命题，只需证明 \(n=1\) 的情形，一般情形由归纳得到。所以可以假设 \(X\) 是某个有限型局部自由 \(\mathcal O_S\)-模 \(E\) 的旗概形。

% 此处构造已核对法文原书第391页。
把 \(S\) 分解成有限多个既开又闭的子概形以后，可以假设 \(E\) 的秩为常数 \(r\)。此时可如下构造概形 \(X\)。令 \(E=F_0\)，\(S=X_0\)，然后依次定义概形 \(X_1,\ldots,X_r\)，并对每个 \(i\) 定义有限型局部自由 \(\mathcal O_{X_i}\)-模 \(F_i\)，方法为
\[
  X_i=\mathbb P(F_{i-1}),\qquad
  F_i=\Ker\bigl(f_i^*(F_{i-1})\longrightarrow \mathcal O_{X_i}(1)\bigr),
\]
其中 \(f_i:X_i\to X_{i-1}\) 是结构态射。因此，只需对射影丛证明命题。视 \(K^\bullet\) 所采用的定义而定，这正是 \hyperlink{sga6.e06.thm.1-1}{1.1} 或 \hyperlink{sga6.e06.thm.2-1}{2.1} 的内容。

\sgaStableTarget{sga6.e06.thm.3-2}\textbf{定理 3.2.} 设 \(S\) 是拟紧概形。由有限型局部自由 \(\mathcal O_S\)-模定义的预 \(\lambda\)-环 \(K^\bullet(S)\)，是一个 \(\lambda\)-环。

同态
\[
  \lambda_t:K^\bullet(S)\longrightarrow 1+K^\bullet(S)[[t]]^+
\]
是可加的，因而为了验证它是 \(\lambda\)-同态，只需在形如 \(\cl(E)\) 的元素上验证普遍关系，其中 \(E\) 是有限型局部自由 \(\mathcal O_S\)-模。因此，要证明：对每一对这样的模 \(E,F\)，仍用 \(E,F\) 表示其类时，关系式 \hyperlink{sga6.e05.eq.2-4-2}{V 2.4.2} 成立：
\[
  \lambda^i(EF)=P_i\bigl(\lambda^1(E),\ldots,\lambda^i(E);\lambda^1(F),\ldots,\lambda^i(F)\bigr),
\]
\[
  \lambda^j(\lambda^i(E))=Q_{i,j}\bigl(\lambda^1(E),\ldots,\lambda^{ij}(E)\bigr).
\]
由于 \(S\) 拟紧，它有一个由两两不交的开集组成的有限覆盖，使 \(E\) 和 \(F\) 在每个开集上都有常秩。由于有限不交并的 \(K^\bullet\) 是相应各 \(K^\bullet\) 的乘积，只需在此覆盖的各成员上验证这些关系；所以可以假设 \(E\) 和 \(F\) 在 \(S\) 上均有常秩。

于是令
\[
  X=D(E)\times_S D(F).
\]
由\hyperlink{sga6.e06.prop.3-1}{命题 3.1}，\(\lambda\)-同态 \(K^\bullet(S)\to K^\bullet(X)\) 为单同态。因此，为了证明关系式 \hyperlink{sga6.e05.eq.2-4-2}{V 2.4.2}，可以在 \(K^\bullet(X)\) 中进行。但在 \(K^\bullet(X)\) 中，元素 \(E\) 和 \(F\) 分裂为可逆层的类之和。由可加性，只需对可逆层的类 \(\xi,\eta\) 证明同样的关系。把关系写成 \hyperlink{sga6.e05.eq.2-4-1}{V 2.4.1} 的形式，就是
\[
  1+\xi\eta t=(1+\xi t)\circ(1+\eta t),
  \qquad
  \lambda_u(1+\xi t)=1+\{1+\xi t\}u.
\]
这些正是定义 \(1+K^\bullet(S)[[t]]^+\) 上的 \(\lambda\)-环结构的关系式 \hyperlink{sga6.e05.eq.2-3-1}{V 2.3.1} 和 \hyperlink{sga6.e05.eq.2-3-3}{V 2.3.3}。

再记下如下更一般的结果。

\sgaStableTarget{sga6.e06.thm.3-3}\textbf{定理 3.3.} 设 \((T,A)\) 是赋交换环拓扑斯，\(K\) 是 \(T\) 上秩有界的局部自由 \(A\)-模的 Grothendieck 群。则带有其典范预 \(\lambda\)-结构的 \(K\) 是 \(\lambda\)-环。

为了证明 \(K\) 是 \(\lambda\)-环，需要证明关系式 \hyperlink{sga6.e05.eq.2-4-1}{V 2.4.1}；由可加性，只需对秩有界的局部自由 \(A\)-模 \(E,F\) 的类证明它们。由于 \(E\) 和 \(F\) 的秩有界，可以把 \(T\) 分解成有限个开集的不交和，使 \(E\) 和 \(F\) 在每个开集上都有常秩。利用 \(K\) 与有限不交和的相容性，并把关系式 \hyperlink{sga6.e05.eq.2-4-1}{V 2.4.1} 写成 \hyperlink{sga6.e05.eq.2-4-2}{V 2.4.2} 的形式，就归结到了常秩模的情形。

设 \(E\) 是秩为 \(n\) 的局部自由 \(A\)-模。有一个 \(\GL(n,A)\)-挠子 \(P\)，使得
\[
  E=P\times_{\GL(n,A)} A^n,
\]
其中通过标准表示把 \(A^n\) 看作 \(\GL(n,A)\)-模。于是可定义从表示环 \(R_{\mathbb Z}(\GL(n))\) 到 \(K\) 的 \(\lambda\)-同态：把 \(\GL(n)\) 在秩为 \(p\) 的自由 \(\mathbb Z\)-模上的表示 \(\rho\) 送到
\[
  P\times_{\GL(n,A)} A^p.
\]
的类。在此同态下，\(E\) 的类是标准表示的像。因此，\(\lambda^i(E)\) 之间的任何多项式关系，都由 \(R_{\mathbb Z}(\GL(n))\) 中标准表示的相应关系推出。

对秩分别为 \(n,p\) 的两个局部自由 \(A\)-模 \(E,F\)，也可作同样的论证。可找到一个 \(\GL(n,A)\times\GL(p,A)\)-挠子 \(P\)，使
\[
 E=P\times_{\GL(n,A)\times\GL(p,A)}A^n,
\]
其中在 \(A^n\) 上，\(\GL(n,A)\) 通过标准表示作用，\(\GL(p,A)\) 通过平凡表示作用；对 \(F\) 则交换两个因子的作用。接着，按照上面相同的方法定义 \(\lambda\)-同态
\[
 R_{\mathbb Z}(\GL(n)\times\GL(p))\longrightarrow K.
\]
于是，\(\lambda^i(E),\lambda^j(F)\) 之间的关系，就归结为 \(R_{\mathbb Z}(\GL(n)\times\GL(p))\) 中相应元素的关系。因此，我们归结到证明：对 \(\Spec\mathbb Z\) 上每个为若干线性群概形的有限乘积的群概形 \(G\)，预 \(\lambda\)-环 \(R_{\mathbb Z}(G)\) 是 \(\lambda\)-环。这可以由 [5] 中对 \(R_{\mathbb Z}(G)\) 的显式计算得到；也可利用 1.14，仿照 \hyperlink{sga6.e06.thm.3-2}{3.2} 进行论证。

\sgaStableTarget{sga6.e06.sec.4}\section*{4. 旗丛的 \(K^\bullet\) 的计算}

\sgaStableTarget{sga6.e06.sec.4.1}\subsection*{4.1}
设 \(S\) 是概形，\(E\) 是秩为 \(n\) 的局部自由 \(\mathcal O_S\)-模。设
\[
  \pi=(p_1,\ldots,p_k),\qquad p_i\geq 1,
  \qquad \sum_{i=1}^k p_i=n.
\]
\(E\) 的 \(\pi\) 型旗函子是 \(S\)-概形范畴上的集合值函子，它把一个 \(S\)-概形 \(f:X\to S\) 送到 \(f^*E\) 的商链
\[
  (L_{k-1},\ldots,L_1),
\]
组成的集合，其中对每个 \(i\)，\(L_i\) 是 \(L_{i+1}\) 的一个局部自由商，秩为 \(p_1+\cdots+p_i\)。此函子由一个 \(S\)-概形 \(D_\pi(E)\) 表示，后者带有 \(E\) 的逆像的典范商序列，各项局部自由，秩分别为 \(p_1,p_1+p_2,\ldots\)。若 \(\pi=(1,n-1)\)，则 \(D_\pi(E)=\mathbb P(E)\)；若 \(\pi=(r,n-r)\)，则 \(D_\pi(E)=G_{r,n}(E)\)，即 \(E\) 的 \((r,n-r)\) 型 Grassmann 概形。若所有 \(p_i\) 均为 \(1\)，则得到 \(E\) 的通常的旗丛，记为 \(D(E)\)。

\sgaStableTarget{sga6.e06.lem.4-2}\textbf{引理 4.2.} 设 \(S\) 是概形，\(E\) 是秩为 \(n\) 的局部自由 \(\mathcal O_S\)-模，并设
\[
  \pi=(p_1,\ldots,p_i,\ldots,p_k),\qquad p_1+\cdots+p_k=n.
\]
设 \(p',p''\geq 1\)，且 \(p'+p''=p_i\)，令
\[
  \pi'=(p_1,\ldots,p_{i-1},p',p'',p_{i+1},\ldots,p_k),
  \qquad \pi''=(p',p'').
\]
设 \(L_1,\ldots,L_{k-1}\) 是 \(E\) 的逆像在 \(D_\pi(E)\) 上的典范商，\(F\) 是 \(L_i\to L_{i-1}\) 的核。则有典范的 \(S\)-同构
\[
  D_{\pi'}(E)\simeq D_{\pi''}(F).
\]
从旗概形的函子性定义立即可得的证明，留给读者。

\sgaStableTarget{sga6.e06.lem.4-3}\textbf{引理 4.3.} 设 \(A\) 是有单位元的交换环，\(n>0\)，\((c_i)_{1\leq i\leq n}\) 是 \(A\) 的元素。设 \(C\) 是由 \(\xi_1,\ldots,\xi_n\) 生成，并满足关系
\[
  \sigma_i(\xi_1,\ldots,\xi_n)=c_i,
  \qquad i=1,\ldots,n,
\]
的 \(A\)-代数，其中 \(\sigma_i\) 表示初等对称函数。设 \(p_1,\ldots,p_k\) 是和为 \(n\) 的整数，均 \(\geq 1\)。对 \(1\leq j\leq k\)，令
\[
  \xi^{(j)}=(\xi^{(j)}_1,\ldots,\xi^{(j)}_{p_j}),
  \qquad
  \xi_i^{(j)}=\xi_{\sum_{\alpha<j}p_\alpha+i}.
\]
令
\[
  c_i^{(j)}=\sigma_i(\xi^{(j)}_1,\ldots,\xi^{(j)}_{p_j}),
  \qquad 1\leq j\leq k,
  \quad 1\leq i\leq p_j,
\]
并设 \(B\) 是 \(C\) 中由各 \(c_i^{(j)}\) 生成的子代数。则：

\begin{enumerate}[label=\textup{\roman*)}]
\item 元素
\[
  (\xi^{(1)})^{\alpha(1)}\cdots(\xi^{(k)})^{\alpha(k)},
  \qquad 0\leq \alpha_i^{(j)}\leq p_j-i,
\]
构成 \(C\) 作为 \(B\)-模的一组基；这里 \(\alpha(j)=(\alpha_1^{(j)},\ldots,\alpha_{p_j}^{(j)})\)。
\item \(A\)-代数 \(B\) 的生成元 \(c_i^{(j)}\) 满足由
\[
% 法文原书第395页的原印公式；下附显式译校注。
  \sum_{i=1}^n c_i
  =
  \prod_{j=1}^{k}\left(\sum_{i=1}^{p_j} c_i^{(j)}\right),
\]
导出的 \(n\) 个等权关系，其中给 \(c_i\) 和 \(c_i^{(j)}\) 赋予权 \(i\)；这 \(n\) 个关系生成各 \(c_i^{(j)}\) 之间的关系理想。\par\smallskip\noindent\textit{译校注：}法文原书第 395 页上式两边均未印常数项。按此处意指的 Chern 类恒等式，应写成 \(1+\sum_i c_i=\prod_j(1+\sum_i c_i^{(j)})\)。正文保留原印式。
\end{enumerate}

\sgaStableTarget{sga6.e06.cor.4-4}\textbf{推论 4.4.} \(C\) 中形如
\[
  \xi_1^{\alpha_1}\cdots \xi_{n-1}^{\alpha_{n-1}},
  \qquad 0\leq \alpha_i\leq n-i,
\]
的元素构成 \(C\) 作为 \(A\)-模的一组基。

引理和推论借自 [2]，第 4 讲，第 4、19 页，其证明参见该处。

\sgaStableTarget{sga6.e06.sec.4.5}\subsection*{4.5}
设 \(S\) 拟紧，\(E\) 是秩为 \(n\) 的局部自由 \(\mathcal O_S\)-模，\(\pi=(p_1,\ldots,p_k)\) 是和为 \(n\) 的整数序列，每项均 \(\geq 1\)，并令 \(X=D_\pi(E)\)。概形 \(X\) 带有典范商 \(L_i\)。令 \(L_0=0\)，\(L_k=f^*E\)，其中 \(f:X\to S\) 是结构态射。对 \(1\leq i\leq k\)，设 \(F_i\) 是 \(L_i\to L_{i-1}\) 的核。设 \(K^\bullet(X)\) 是 \(X\) 上有限型局部自由层、或完美复形的 Grothendieck 群。用 \(F_j\) 表示 \(F_j\) 的类，用 \(\lambda_i^{(j)}\) 表示 \(\lambda^i(F_j)\) 的类。当 \(K^\bullet(X)\) 由完美复形定义时，给级数 \(\lambda_t(E)\) 和 \(\lambda_t(F_j)\) 赋予同样显然的意义，其中 \(E\) 表示 \(E\) 的类。最后，
\[
  \sum_{j=1}^{k}F_j=E.
\]

\sgaStableTarget{sga6.e06.thm.4-6}\textbf{定理 4.6.} 在 \hyperlink{sga6.e06.sec.4.5}{4.5} 的假设和记号下，\(K^\bullet(S)\)-代数 \(K^\bullet(X)\) 由 \(\lambda_i^{(j)}\) 生成，而这些 \(\lambda_i^{(j)}\) 之间的关系理想由以下等式导出的关系生成：
\[
  \lambda_t(E)=\prod_{j=1}^{k}\lambda_t(F_j).
\]
等价地，\(K^\bullet(X)\) 作为 \(K^\bullet(S)\)-\(\lambda\)-代数（参见 V 4.6），由 \(F_1,\ldots,F_k\) 生成，并满足关系
\[
  \text{(i)}\quad \lambda^p(F_i)=0\quad(p>p_i),
  \qquad
  \text{(ii)}\quad \sum_i F_i=E.
\]

我们先处理所有 \(p_i\) 都等于 \(1\) 的情形。此时定理可表述如下。

\sgaStableTarget{sga6.e06.cor.4-7}\textbf{推论 4.7.} 设 \(S\) 拟紧，\(E\) 是秩为 \(n\) 的局部自由 \(\mathcal O_S\)-模，\(X=D(E)\) 是 \(E\) 的旗丛，\(\xi_1,\ldots,\xi_n\) 是可逆层 \(F_1,\ldots,F_n\) 的类，采用 4.5 的记号。\footnote{译校注：原书第 397 页此处将 4.5 误排为 3.5；同页证明中的“推论 3.4”应为推论 4.4，第 398 页“3.5 和 3.7 的记号”应为“4.5 和 4.7 的记号”。}则 \(K^\bullet(X)\) 是由 \(\xi_i\) 生成，并满足关系
\[
  \sigma_i(\xi_1,\ldots,\xi_n)=\lambda^i(E).
\]
的 \(K^\bullet(S)\)-代数。
对 \(n\) 作归纳。当 \(n=1\) 时，断言显然。假设已知 \(n-1\) 时成立。考虑 \(\mathbb P(E)\) 以及其上的正合列
\[
  0\longrightarrow F\longrightarrow g^*E\longrightarrow \mathcal O_{\mathbb P(E)}(1)\longrightarrow 0,
\]
其中 \(g\) 是 \(\mathbb P(E)\) 的结构态射。由\hyperlink{sga6.e06.lem.4-2}{引理 4.2}，\(X\) 与 \(\mathbb P(E)\) 上通常的旗丛 \(D(F)\) 等同，而且在此等同下，典范层彼此对应。图
\[
\begin{tikzcd}[column sep=large,row sep=large]
D(E)=D(F) \arrow[r,"h"] \arrow[dr,"f"'] & \mathbb P(E) \arrow[d,"g"]\\
& S
\end{tikzcd}
\]
表明 \(K^\bullet(S)\to K^\bullet(X)\) 通过 \(K^\bullet(\mathbb P(E))\) 分解。对 \(D(F)\) 应用归纳假设，再由\hyperlink{sga6.e06.cor.4-4}{推论 4.4}，\(K^\bullet(X)\) 在 \(K^\bullet(\mathbb P(E))\) 上有一组基，由
\[
  \xi_2^{\alpha_2}\cdots \xi_{n-1}^{\alpha_{n-1}},
  \qquad 0\leq \alpha_i\leq n-i.
\]
组成。此外，\(K^\bullet(\mathbb P(E))\) 在 \(K^\bullet(S)\) 上有基 \(1,\xi_1,\ldots,\xi_1^{n-1}\)。因此，\(K^\bullet(X)\) 在 \(K^\bullet(S)\) 上有基
\[
  \xi_1^{\alpha_1}\cdots \xi_{n-1}^{\alpha_{n-1}},
  \qquad 0\leq \alpha_i\leq n-i.
\]
另外，\(\lambda_t\) 的可加性给出 \(\prod_i(1+\xi_i t)=\lambda_t(E)\)，故所列的关系均成立。若 \(C\) 是由 \(\xi_i\) 生成并满足这些关系的 \(K^\bullet(S)\)-代数，则有满同态 \(C\to K^\bullet(X)\)。由\hyperlink{sga6.e06.cor.4-4}{推论 4.4}，此同态把 \(C\) 的一组基送到 \(K^\bullet(X)\) 的一组基，因而它是同构。

对一般的 \(\pi\) 型旗丛，采用 4.5 和 4.7 的记号，从 \(D_\pi(E)\) 出发依次构造各 \(F_j\) 的旗丛，即得到通常的旗丛 \(D(E)\)。因此，\(K^\bullet(D(E))\) 作为 \(K^\bullet(X)\)-模，有一组基由 4.3.i) 中的元素组成，其中 \(\xi_i^{(j)}\) 是在 \(D(E)\) 上将各 \(F_j\) 的逆像分裂时出现的、类似于各 \(F_j\) 的层的类。因此，\(K^\bullet(X)\) 是 \(K^\bullet(D(E))\) 的子环。此外，由可加性，
\[
  \lambda_i^{(j)}=\sigma_i(\xi_1^{(j)},\ldots,\xi_{p_j}^{(j)}),
\]
因此，\(K^\bullet(X)\) 包含 \(K^\bullet(D(E))\) 中由各 \(\lambda_i^{(j)}\) 生成的 \(K^\bullet(S)\)-子代数。若 \(B\) 是此子代数，则由 4.3.i)，同一族元素对 \(B\)-模结构及 \(K^\bullet(X)\)-模结构都构成一组基。因此 \(K^\bullet(X)=B\)，这就证明了生成性。各 \(\lambda_i^{(j)}\) 之间的关系理想由 4.3.ii) 给出，此时利用了 \hyperlink{sga6.e06.cor.4-7}{4.7} 确定的 \(K^\bullet(D(E))\) 在 \(K^\bullet(S)\) 上的结构，以及 \(\xi_i^{(j)}\) 恰好是 \(E\) 的逆像在 \(D(E)\) 上分裂中出现的各类这一事实。这就证明了 4.6 的第一种表述。

对于第二种表述，注意到 \(K^\bullet(S)\)-\(\lambda\)-代数 \(K^\bullet(X)\) 的元素 \(F_i\) 满足 (i) 和 (ii)。若 \(A\) 是由满足 (i) 和 (ii) 的元素 \(F_i'\) 生成的 \(K^\bullet(S)\)-\(\lambda\)-代数，则有 \(K^\bullet(S)\)-\(\lambda\)-同态 \(A\to K^\bullet(X)\)，把 \(F_i'\) 送到 \(F_i\)。由 \hyperlink{sga6.e06.thm.4-6}{4.6}，此同态是满同态。由于 (i) 和 (ii) 在 \(K^\bullet(S)[F_1,\ldots,F_k]\) 中生成的 \(\lambda\)-理想包含第一种表述中的关系生成的通常理想，该同态也是单同态。因此它是同构。

\sgaStableTarget{sga6.e06.cor.4-8}\textbf{推论 4.8.} 在 \hyperlink{sga6.e06.sec.4.5}{4.5} 的假设和记号下，令
\[
  c_i^{(j)}=\gamma^i(F_j-p_j),
  \qquad 1\leq i\leq p_j.
\]
则 \(K^\bullet(S)\)-代数 \(K^\bullet(X)\) 由各 \(c_i^{(j)}\) 生成，而它们之间的关系理想由以下等式导出的关系生成：
\[
  \gamma_t(E-n)=\prod_{j=1}^{k}\gamma_t(F_j-p_j).
\]
这只是 4.6 的另一种表述。事实上，各 \(c_i^{(j)}\) 可以表示为各 \(\lambda_i^{(j)}\) 的函数，反之亦然。此外，关系 \(\lambda_t(E)=\prod\lambda_t(F_j)\) 等价于 \(\gamma_t(E-n)=\prod\gamma_t(F_j-p_j)\)。

\sgaStableTarget{sga6.e06.sec.4.9}\subsection*{4.9}
设 \(S\) 是概形。对 \(r\ge n\)，用 \(G_S^{n,r}\) 表示 \(\mathcal O_S^r\) 的 \((n,r-n)\) 型 Grassmann 概形。使 \(r\) 变化时，各 \(G_S^{n,r}\) 构成归纳系，其中态射
\[
  G_S^{n,r}\longrightarrow G_S^{n,r+1}
\]
由把 \(\mathcal O_{G_S^{n,r}}^r\) 看作 \(\mathcal O_{G_S^{n,r}}^{r+1}\) 除以第 \((r+1)\) 个因子所得的商来定义。因此，在此态射系中，典范商通过逆像彼此对应。这样定义的归纳系记为 \(G_S^{n,\infty}\)。

我们用公式
\sgaStableTarget{sga6.e06.eq.4-9-1}\begin{equation}\tag{4.9.1}
  K^\bullet(G_S^{n,\infty})
  = \varprojlim_r K^\bullet(G_S^{n,r}) .
\end{equation}
定义其 \(K^\bullet\)。设 \(F_1,F_2\) 是 \(G_S^{n,r}\) 上的典范层 \hyperlink{sga6.e06.sec.4.5}{(4.5)}，也用 \(F_1,F_2\) 表示其在 \(K^\bullet(G_S^{n,r})\) 中的类。令
\[
  c_i=\gamma^i(F_1-n).
\]
当 \(r\) 变化时，各 \(c_i\) 在过渡态射下彼此对应，因而定义 \(K^\bullet(G_S^{n,\infty})\) 中的元素，仍记为 \(c_i\)。

\sgaStableTarget{sga6.e06.prop.4-10}\textbf{命题 4.10.} 设 \(S\) 是拟紧概形。则在 \hyperlink{sga6.e06.sec.4.9}{4.9} 的记号下，有同构
\[
  K^\bullet(G_S^{n,\infty})\simeq K^\bullet(S)[[c_1,\ldots,c_n]] .
\]

由 \hyperlink{sga6.e06.cor.4-8}{4.8}，\(K^\bullet(G_S^{n,r})\) 是由
\[
  c_i=\gamma^i(F_1-n),\quad 1\le i\le n,
\]
以及
\[
  \gamma^j(F_2-r+n),\quad 1\le j\le r-n,
\]
生成的 \(K^\bullet(S)\)-代数，其关系由
\[
  \gamma_t(F_1-n)\cdot\gamma_t(F_2-r+n)
  = \gamma_t(r-r)=1,
\]
导出，或等价地由
\[
  \gamma_t(F_2-r+n)=1/\gamma_t(F_1-n).
\]
导出。此关系可以把各 \(\gamma^j(F_2-r+n)\) 表示为各 \(c_i\) 的式子；写出 \(1/\gamma_t(F_1-n)\) 的所有次数 \(>r-n\) 的系数为零，便给出各 \(c_i\) 之间的一系列等权关系（赋予 \(c_i\) 权 \(i\)）；而且这些关系生成各 \(c_i\) 的关系理想。因此可写为
\[
  K^\bullet(G_S^{n,r})=K^\bullet(S)[c_1,\ldots,c_n]/I_r,
\]
其中 \(I_r\) 是 \(K^\bullet(S)[c_1,\ldots,c_n]\) 的次数 \(>r-n\) 的齐次理想。由\hyperlink{sga6.e06.eq.4-9-1}{定义 (4.9.1)} 可知，\(K^\bullet(G_S^{n,\infty})\) 是 \(K^\bullet(S)[c_1,\ldots,c_n]\) 对各 \(I_r\) 所定义拓扑的分离完备化。此拓扑比由权定义的拓扑更细；我们将证明二者等价，这就推出命题。

令 \(A=K^\bullet(S)[c_1,\ldots,c_n]/I_r\)，设 \(C\) 是由生成元 \(x_1,\ldots,x_n\) 生成、满足关系
\[
  \sigma_j(x_1,\ldots,x_n)=c_j,
\]
的 \(A\)-代数，其中 \(\sigma_j\) 表示初等对称函数。由 \hyperlink{sga6.e06.cor.4-4}{4.4}，典范同态 \(A\to C\) 为单同态。得到
\[
  \gamma_t(F_1-n)
  =1+c_1t+\cdots+c_nt^n
  =\prod_{i=1}^n(1+x_it).
\]
由于 \(1/\gamma_t(F_1-n)\) 是多项式，对每个 \(i\)，\(1/(1+x_it)\) 也是多项式。因此，各 \(x_i\) 是幂零元，且存在整数 \(N\)，使每个以 \(K^\bullet(S)\) 为系数、关于各 \(x_i\) 的次数 \(>N\) 的多项式均为零。因此，\(A\) 中每个权 \(>N\) 的多项式均为零，即 \(K^\bullet(S)[c_1,\ldots,c_n]\) 中每个权 \(>N\) 的多项式都属于 \(I_r\)。所以所考虑的两种拓扑等价，证明完成。

\sgaStableTarget{sga6.e06.sec.5}\section*{5. 对射影丛的应用；\(f_*\) 的研究}

\sgaStableTarget{sga6.e06.sec.5.1}\subsection*{5.1}
设 \(S\) 拟紧，\(E\) 局部自由且秩为 \(r+1\)，并令 \(X=\mathbb P(E)\)，其结构态射为 \(f:X\to S\)。在 \(X\) 上有典范正合列
\[
  0\longrightarrow F\longrightarrow f^*E\longrightarrow \mathcal O_X(1)\longrightarrow 0.
\]
本节的 \(K^\bullet\) 均用有限型局部自由层定义。用 \(E,F,\xi\) 表示相应的类，并令 \(x=1-\xi^{-1}\)，\(x'=\xi-1=\xi x\)，用 \(\bar x,\bar x'\) 表示其在 \(\Gr^1K^\bullet(X)\) 中的类。由于 \(\xi-1\) 和 \(x\) 都属于 \(\Fil^1K^\bullet(X)\)，有 \((\xi-1)x\in\Fil^2K^\bullet(X)\)，即 \(x'=\xi x\equiv x\pmod{\Fil^2K^\bullet(X)}\)，从而 \(\bar x'=\bar x\)。

\sgaStableTarget{sga6.e06.sec.5.2}\subsection*{5.2}
对 \(X\) 上的有限型局部自由层 \(G\)，高阶正像 \(R^if_*G\) 不一定是 \(S\) 上的有限型局部自由层。因此，不能像用完美复形定义 \(K^\bullet\) 时那样，用高阶正像的交错和定义正像 \(f_*:K^\bullet(X)\to K^\bullet(S)\)。\hyperlink{sga6.e06.thm.1-1}{定理 1.1} 使我们可以如下绕过这个困难：每个 \(y\in K^\bullet(X)\) 可以唯一地表示成
\sgaStableTarget{sga6.e06.eq.5-2-1}\begin{equation}\tag{5.2.1}
   y=a_0+a_1\xi^{-1}+\cdots+a_r\xi^{-r},
  \qquad a_i\in K^\bullet(S),
\end{equation}
于是令
\sgaStableTarget{sga6.e06.eq.5-2-2}\begin{equation}\tag{5.2.2}
  f_*(y)=a_0.
\end{equation}
此加法同态满足投影公式。此外，
\sgaStableTarget{sga6.e06.eq.5-2-3}\begin{equation}\tag{5.2.3}
\begin{aligned}
  f_*(1)&=1,\qquad f_*(\xi^n)=0\quad(-r\leq n<0),\\
  f_*(\xi^n)&=s_n\quad(n\geq0),
\end{aligned}
\end{equation}
其中 \(s_n\) 是 \(E\) 的各根的第 \(n\) 个完全对称函数。为验证最后一组关系，对 \(n\) 作归纳。假设对 \(p\leq n-1\) 已有 \(f_*(\xi^p)=s_p\)。把关系式 \hyperlink{sga6.e06.eq.1-9-1}{(1.9.1)} 乘以 \(\xi\) 的适当幂，得到 \(\sum_{i=0}^{r+1}(-1)^i\lambda^i(E)\xi^{n-i}=0\)。应用 \(f_*\) 后得到 \(\sum_{i=0}^{r+1}(-1)^i\lambda^i(E)f_*(\xi^{n-i})=0\)。由关系式 \hyperlink{sga6.e06.eq.1-11-3}{(1.11.3)} 和归纳假设，即得结论。

\sgaStableTarget{sga6.e06.prop.5-3}\textbf{命题 5.3.} 在 5.1 的假设下，\footnote{译校注：原书第 403 页将这里的 5.1 误排为 4.1。}对每个 \(k\geq 0\)，有
\[
  \Fil^k K^\bullet(X)
  =
  \Fil^k K^\bullet(S)
  +\Fil^{k-1}K^\bullet(S)x+\cdots+\Fil^{k-r}K^\bullet(S)x^r,
\]
这是 \(K^\bullet(S)\)-模的直和分解，其中滤过是增广 \(\lambda\)-环的典范滤过。

首先注意到，\(x\) 的增广为零。另外，
\[
 \gamma_t(x)=\lambda_{t/(1-t)}(1-\xi^{-1})
 =\left(1+\frac{t}{1-t}\right)
  \left(1+\xi^{-1}\frac{t}{1-t}\right)^{-1}
 =\frac1{1-xt}.
\]
因此，对每个 \(i\)，有 \(\gamma^i(x)=x^i\)。所以，\(K^\bullet(X)\) 是由生成元 \(x\) 及关系 \(\gamma^i(x)=x^i\) 生成的 \(K^\bullet(S)\)-\(\lambda\)-代数的一个商，而取商所用的理想由关系 (1.13.1) 生成。于是可应用 V 4.12 和 V 4.13.ii)，只须证明，对每个 \(n\geq0\)，
\[
 x^n\in\Fil^nK^\bullet(S)+\Fil^{n-1}K^\bullet(S)x
       +\cdots+\Fil^{n-r}K^\bullet(S)x^r.
\]
当 \(n\leq r\) 时显然成立；以后对 \(n\) 归纳。关系 (1.13.1) 给出
\[
 x^n=-\sum_{i=0}^r\gamma^{r+1-i}(\check E-r-1)x^{n-r-1+i},
 \qquad n\geq r+1.
\]
对各 \(x^{n-r-1+i}\) 应用归纳假设，即得所需结论。由于 \(1,\ldots,x^r\) 在 \(K^\bullet(S)\) 上线性无关，此和为直和。

\sgaStableTarget{sga6.e06.cor.5-4}\textbf{推论 5.4.} \(K^\bullet(X)\) 在 \(K^\bullet(S)\) 上诱导的滤过，等于 \(K^\bullet(S)\) 原有的滤过。

\sgaStableTarget{sga6.e06.cor.5-5}\textbf{推论 5.5.} 若 \(X\) 是若干局部自由层的旗丛在 \(S\) 上的乘积，则
\[
  \Gr K^\bullet(S)\longrightarrow \Gr K^\bullet(X)
\]
是单同态。

按照 3.1 的论证，可归结到 \(X\) 是 \(S\) 上一个射影丛的情形；此情形直接由 5.4 得到。

\sgaStableTarget{sga6.e06.cor.5-6}\textbf{推论 5.6.} 使用上述射影丛的记号，\(\Gr K^\bullet(X)\) 是以 \(1,x,\ldots,x^r\) 为基的自由 \(\Gr K^\bullet(S)\)-模，并且
\[
  \Gr^k K^\bullet(X)
  =\Gr^kK^\bullet(S)\oplus\Gr^{k-1}K^\bullet(S)x\oplus\cdots\oplus\Gr^{k-r}K^\bullet(S)x^r.
\]
元素 \(1,x,\ldots,x^{r+1}\) 满足
\[
  \sum_{i=0}^{r+1}(-1)^i c_i(E)x^{r+1-i}=0,
\]
其中 \(c_i(E)\) 是 \(E\) 的第 \(i\) 个 Chern 类（V 6.7）。这是 1.1、5.3 及 (1.13.1) 的直接推论。

\sgaStableTarget{sga6.e06.cor.5-7}\textbf{推论 5.7.} 对 \hyperlink{sga6.e06.sec.4.5}{4.5} 中的旗丛，\(\Gr K^\bullet(S)\)-代数 \(\Gr K^\bullet(X)\) 由各 \(\gamma^i(F_\alpha-p_\alpha)\) 在 \(\Gr^iK^\bullet(X)\) 中的类 \(\bar\gamma_i^{(\alpha)}\) 生成，其关系由
\[
   c_t(E)=\prod_{\alpha}c_t(F_\alpha),
\]
导出，其中 \(c_t\) 表示总 Chern 类。照搬 4.6 的证明并利用 5.6，即可得到此结果（比较 [2]，第 4 讲）。

\sgaStableTarget{sga6.e06.cor.5-8}\textbf{推论 5.8.} 对每个 \(k\)，
\[
   f_*(\Fil^k K^\bullet(X))\subset \Fil^{k-r} K^\bullet(S).
\]
这由\hyperlink{sga6.e06.prop.5-3}{命题 5.3} 和投影公式立即得到。

\sgaStableTarget{sga6.e06.prop.5-9}\textbf{命题 5.9.} 使用 \hyperlink{sga6.e06.sec.5.1}{5.1} 的记号和假设，对每个 \(i\geq1\)，有
\[
   f_*(\lambda^i(F))=0,
  \qquad
   f_*(\gamma^i(F))=0.
\]
特别地，令 \(\lambda_{-1}(F)=\sum_{i=0}^r(-1)^i\lambda^i(F)\)，则得到
\[
  f_*(\lambda_{-1}(F))=1.
\]
按系数定义映射
\[
  \xi_t:1+K^\bullet(X)[[t]]^+\longrightarrow 1+K^\bullet(S)[[t]]^+
\]
如下：
\[
  \xi_t\left(\sum a_i t^i\right)=\sum f_*(a_i)t^i,
\]
虽然它一般不是群同态，但对 \(\varphi\in1+K^\bullet(X)[[t]]^+\) 和 \(\psi\in1+K^\bullet(S)[[t]]^+\)，它满足
\[
   \xi_t\bigl(\varphi\,f^*(\psi)\bigr)=\psi\,\xi_t(\varphi).
\]
事实上，这逐个系数成立，直接来自投影公式。由 \(E=F+\xi\)，得到 \(\xi_t(\lambda_t(F))=\xi_t(\lambda_t(E)\lambda_t(-\xi))=\lambda_t(E)\xi_t(\lambda_t(-\xi))\)。由 (5.2.3)，\footnote{译校注：原书第 406 页此处将 (5.2.3) 误排为 (4.2.3)。}\(\xi_t(\lambda_t(-\xi))=\sum_{i\geq0}(-t)^if_*(\xi^i)=\sigma_{-t}(E)\)。利用关系式 (1.12.1)，最终得到 \(\xi_t(\lambda_t(F))=1\)，由此即得所需各关系。

\sgaStableTarget{sga6.e06.prop.5-10}\textbf{命题 5.10.} 设 \(z\in K^\bullet(X)\)。条件
\[
  z(1-\xi)=0
\]
等价于存在元素 \(y\in K^\bullet(S)\)，使得
\[
  \text{(i)}\quad z=y\lambda_{-1}(F),
  \qquad
  \text{(ii)}\quad y\lambda_{-1}(E)=0.
\]
充分性来自
\[
  \lambda_{-1}(E)=\lambda_{-1}(F)\lambda_{-1}(\xi)=\lambda_{-1}(F)(1-\xi).
\]
设 \(z(1-\xi)=0\)，即 \(zx'=0\)。由注 1.13，存在唯一的 \(a_i\in K^\bullet(S)\)，使
\[
 z=\sum_{i=0}^r a_ix'{}^i.
\]
另一方面，\(x'{}^i\) 满足关系
\[
 \sum_{i=0}^r c_ix'{}^i+x'{}^{r+1}=0,
 \qquad c_i=(-1)^{r+1+i}\gamma^{r+1-i}(E-r-1).
\]
由 \(zx'=0\)，得到
\[
 a_rx'{}^{r+1}=-\sum_{i=1}^r a_{i-1}x'{}^i
             =-\sum_{i=0}^r a_rc_ix'{}^i.
\]
利用 \(1,\ldots,x'{}^r\) 的线性无关性，得
\[
 a_rc_0=0,\qquad a_i=a_rc_{i+1}\quad(0\leq i<r).
\]
因此
\[
 z=a_r\left(\sum_{i=0}^{r-1}c_{i+1}x'{}^i+x'{}^r\right).
\]
而
\[
\begin{aligned}
 \sum_{i=0}^{r-1}c_{i+1}x'{}^i+x'{}^r
 &=(-1)^r\left(\sum_{i=0}^{r-1}
       \gamma^{r-i}(E-r-1)(-1)^ix'{}^i+(-1)^rx'{}^r\right)\\
 &=(-1)^r\sum_{i=0}^r\gamma^{r-i}(E-r-1)\gamma^i(-x')\\
 &=(-1)^r\gamma^r(E-r-1-x')\\
 &=(-1)^r\gamma^r(F-r)\\
 &=(-1)^r\lambda^r(F-1)=\lambda_{-1}(F).
\end{aligned}
\]
令 \(y=a_r\)，即得 \(z=y\lambda_{-1}(F)\)。此外，\(yc_0=0\)，而
\[
 c_0=(-1)^{r+1}\gamma^{r+1}(E-r-1)=\lambda_{-1}(E),
\]
故 \(y\lambda_{-1}(E)=0\)。

\sgaStableTarget{sga6.e06.prop.5-11}\textbf{命题 5.11.} 设 \(z\in \Gr K^\bullet(X)\)。条件 \(zx=0\) 等价于存在 \(y\in \Gr K^\bullet(S)\)，使得
\[
  \text{(i)}\quad z=y\,c_{\mathrm{top}}(F),
  \qquad
  \text{(ii)}\quad y\,c_{\mathrm{top}}(E)=0.
\]
过渡到相伴分次环并利用 5.6 后，证明与 \hyperlink{sga6.e06.prop.5-10}{5.10} 相同。

\sgaStableTarget{sga6.e06.sec.6}\section*{6. \(X\) 有丰沛层时 \(K^\bullet(X)\) 的滤过的研究}

设 \(X\) 是有丰沛层的概形。则 \(K^\bullet(X)\) 的两种定义，即用有限型局部自由层的定义和用完美复形的定义，彼此相同。设
\[
  \varepsilon:K^\bullet(X)\longrightarrow H^0(X,\mathbb Z)
\]
为增广同态。

\sgaStableTarget{sga6.e06.prop.6-1}\textbf{命题 6.1.} 若 \(X\) 拟紧且有丰沛层，则 \(\Fil^1 K^\bullet(X)\) 是幂零元理想。

\(\Fil^1 K^\bullet(X)\) 由形如 \(E-\varepsilon(E)\) 的元素生成，其中 \(E\) 是局部自由层的类，所以只需证明这样的元素幂零。\(X\) 的拟紧性以及 \(K^\bullet\) 与有限直和的相容性，使我们可以归结到 \(E\) 的秩为常数的情形。最后过渡到 \(E\) 的旗丛，利用 \hyperlink{sga6.e06.prop.3-1}{3.1} 给出的单同态，并保留丰沛层的存在性，便可假设 \(E\) 可逆。设 \(E^*\) 是 \(E\) 的对偶。对充分大的 \(n\)，\(E^*(n)\) 由全局截面生成，故可找到满同态 \(\mathcal O_X^m\to E^*(n)\)；与 \(E(-n)\) 张量后，得到满同态 \((E(-n))^m\to\mathcal O_X\)。此同态的核局部自由且秩为 \(m-1\)。用 \(L\) 表示丰沛层的类，便有
\[
   \lambda^m(mEL^{-n}-1)=0.
\]
% 恢复法文原书第 408 页的计算；英文底本标记 SGA-SRC-001543。
而 \(\lambda^m(mEL^{-n}-1)=(-1)^m\lambda_{-1}(mEL^{-n})=(-1)^m\lambda_{-1}(EL^{-n})^m=(-1)^m(1-EL^{-n})^m\)。所以，对每个可逆层 \(E\)，当 \(n\) 足够大时，元素 \(1-EL^{-n}\) 幂零；乘以 \(L^n\) 后得到的 \(L^n-E\) 也幂零。现在可写成 \(E-1=(L^n-1)-(L^n-E)\)。由前面的论证，当 \(n\) 足够大时，\(L^n-E\) 和 \(L^n-1\) 都幂零，因而 \(E-1\) 幂零。

\sgaStableTarget{sga6.e06.prop.6-2}\textbf{推论 6.2.} 在命题 6.1 的假设下，设 \(R\) 是 \(K^\bullet(X)\) 的 \(\lambda\)-子环，由有限多个层的类生成。则 \(R\) 的 \(\lambda\)-环滤过是离散的，且 \(R\cap\Fil^1K^\bullet(X)\) 幂零。此外，\(R\) 在 \(R\cap H^0(X,\mathbb Z)\) 上是有限模（V 3.9.6）。

由 \(X\) 拟紧，可找到由有限个两两不交开集组成的覆盖，使生成 \(R\) 的各层在每个开集上均为常秩。通常的论证于是将问题归结到这些层在 \(X\) 上均为常秩的情形；此时 \(R\) 的增广取值于 \(\mathbb Z\)。

设这些层为 \(E_i\)（\(1\leq i\leq k\)），其类仍记为 \(E_i\)，令 \(x_i=E_i-\varepsilon(E_i)\)，\(r_i=\varepsilon(E_i)\)。则 \(R\) 是由各 \(x_i\) 及关系
\(\gamma^j(x_i)=0\)（\(j>r_i\)）生成的 \(\lambda\)-环的一个商，亦即
\[
 \mathbb Z[x_i,\gamma^j(x_i)]_{\,1\leq i\leq k,\ 2\leq j\leq r_i}
\]
的一个商；所以 \(R\) 诺特。由命题 6.1，各 \(x_i,\gamma^j(x_i)\) 均幂零。因此 \(R\) 的增广理想幂零，而此理想正是 \(R\cap\Fil^1K^\bullet(X)\)。

\(\Fil^nR\) 的每个元素都能表示成关于各 \(x_i,\gamma^j(x_i)\) 的多项式，其权小于 \(n\) 的齐次分量均为零。设 \(m\) 是满足某个 \(\gamma^m(x_i)\ne0\) 的最大整数，取 \(p\) 使所有 \(x_i,\gamma^j(x_i)\) 的 \(p\) 次幂均为零。当 \(n>kp(1+\cdots+m)\) 时，显然有 \(\Fil^nR=0\)；故滤过离散。

滤过既然离散，关于各 \(x_i,\gamma^j(x_i)\) 的非零单项式只有有限多个。这些单项式作为 \(\mathbb Z\)-模生成 \(R\)，证明完成。

\sgaStableTarget{sga6.e06.cor.6-3}\textbf{推论 6.3.} 在命题 6.1 的假设下，设 \(A\) 是交换环。\(K^\bullet(X)\otimes_{\mathbb Z}A\) 的元素 \(y\) 可逆，当且仅当其增广可逆。特别地，每个处处非零的有限型局部自由层的类在 \(K^\bullet(X)\otimes_{\mathbb Z}\mathbb Q\) 中可逆。

事实上，增广理想是幂零元理想。

\sgaStableTarget{sga6.e06.cor.6-4}\textbf{推论 6.4.} 设 \(X\) 拟紧且有丰沛层，\(f:X'\to X\) 固有且完美。再假设 \(f\) 射影或 \(X\) 诺特，使得有限性定理适用。若存在元素 \(x'\in K^\bullet(X')\otimes A\)，使 \(f_*(x')\) 的增广可逆，则
\[
  f^*:K^\bullet(X)\otimes A\longrightarrow K^\bullet(X')\otimes A
\]
是单同态。
事实上，由投影公式，
\[
  f_*(f^*y\cdot x')=y\cdot f_*(x'),
\]
而 \(f_*(x')\) 的增广可逆，故由 6.3，它可逆。因此复合同态\[K^\bullet(X)\otimes A\xrightarrow{f^*}K^\bullet(X')\otimes A\xrightarrow{\,\cdot x'\,}K^\bullet(X')\otimes A\xrightarrow{f_*}K^\bullet(X)\otimes A\]是乘以可逆元 \(f_*(x')\) 的同态，从而是自同构；所以第一个同态 \(f^*\) 为单同态。

\sgaStableTarget{sga6.e06.sec.6.5}\subsection*{6.5. 余维数滤过}
现在假设 \(X\) 诺特，维数为 \(d\)，且有丰沛层。如下定义滤过 \(K^\bullet(X)^i\)：元素 \(x\in K^\bullet(X)\) 属于 \(K^\bullet(X)^i\)，当且仅当对每个闭子概形 \(Y\subset X\)，都可用一个完美复形 \(M^\bullet(Y)\) 表示，使得
\begin{enumerate}[label=\textup{\alph*)}]
\item \(\cl(M^\bullet)=x\);
\item \(\codim(\Supp H^*(M^\bullet(Y))\cap Y,Y)\geq i\)，其中 \(H^*(M^\bullet(Y))=\bigoplus_jH^j(M^\bullet(Y))\)。
\end{enumerate}
这些 \(K^\bullet(X)^i\) 是加法子群。事实上，若 \(x,x'\) 属于此子群，对给定闭集 \(Y\) 取相应复形 \(M^\bullet,M'{}^\bullet\)，则 \(x+x'\) 是 \(M^\bullet\oplus M'{}^\bullet\) 的类，其总上同调的支集是两个总上同调的支集之并。有限个闭集的并的余维数等于它们的余维数的下确界，故 \(x+x'\) 仍属于此子群。

\sgaStableTarget{sga6.e06.prop.6-6}\textbf{命题 6.6.} \hyperlink{sga6.e06.sec.6.5}{6.5} 的滤过\footnote{译校注：原书第 411 页将此处的 6.5 误排为 5.5。}是环滤过，且 \(K^\bullet(X)^{d+1}=0\)。

设 \(x\in K^\bullet(X)^i\)，\(x'\in K^\bullet(X)^j\)，\(Y\) 是 \(X\) 的闭集。先取代表 \(x\) 的完美复形 \(M^\bullet\)，使 \(\codim(\Supp H^*(M^\bullet)\cap Y,Y)\geq i\)；再取代表 \(x'\) 的完美复形 \(M'{}^\bullet\)，使 \(\codim(\Supp H^*(M^\bullet)\cap\Supp H^*(M'{}^\bullet)\cap Y,\Supp H^*(M^\bullet)\cap Y)\geq j\)。则 \(xx'\) 是 \(M^\bullet\otimes^{\mathbf L}M'{}^\bullet\) 的类，且由 EGA \(0_{\mathrm{IV}}\), 14.2.2.3，两个支集与 \(Y\) 的交在 \(Y\) 中的余维数至少为 \(i+j\)。又有
\[
 \Supp H^*(M^\bullet\otimes^{\mathbf L}M'{}^\bullet)\subseteq\Supp H^*(M^\bullet)\cap\Supp H^*(M'{}^\bullet).
\]
因此乘法与此滤过相容。又因大于 \(d\) 的余维数所对应的支集为空，\(K^\bullet(X)^{d+1}=0\)。

\sgaStableTarget{sga6.e06.lem.6-7}\textbf{引理 6.7.} 设 \(X=\Spec A\) 仿射诺特，\(x_1,\ldots,x_n\) 是 \(X\) 的有限多个点，\(E\) 可逆。则存在截面 \(s\in \Gamma(X,E)\)，使它在每个 \(x_i\) 处都是 \(E_{x_i}\) 的一组基。

令 \(M=\Gamma(X,E)\)。对每个 \(i\)，取 \(X\) 的一个闭点 \(x_i'\)，使 \(x_i'\in\overline{\{x_i\}}\)。各 \(x_i'\) 对应于 \(A\) 的极大理想 \(\mathfrak m_i\)；必要时改变指标集，使它们两两不同，每个理想只出现一次。由于 \(E\) 可逆，对每个 \(i\)，有
\[
 M_{\mathfrak m_i}\simeq A_{\mathfrak m_i},
 \qquad M/\mathfrak m_iM\simeq A/\mathfrak m_i.
\]
由此得到
\[
 \prod_iM/\mathfrak m_iM\simeq\prod_i A/\mathfrak m_i,
 \qquad
 M/(\textstyle\prod_i\mathfrak m_i)M
 \simeq A/\textstyle\prod_i\mathfrak m_i.
\]
提升最后这个同构，得到一个同态 \(A\to M\)，其局部化 \(A_{\mathfrak m_i}\to M_{\mathfrak m_i}\) 对每个 \(i\) 模 \(\mathfrak m_i\) 都为同构，因而由 Nakayama 引理，它本身就是同构。\(1\) 的像便给出 \(E\) 的一个截面 \(s\)，在每个 \(x_i'\) 处为一组基。由于 \(s\) 成为 \(E\) 的基的点组成 \(X\) 的一个开集，\(s\) 在每个 \(x_i\) 处也为 \(E_{x_i}\) 的一组基。

\sgaStableTarget{sga6.e06.lem.6-8}\textbf{引理 6.8.} 设 \(X\) 诺特且有丰沛层 \(L\)，\(E\) 可逆，\(x_1,\ldots,x_n\) 是 \(X\) 的点。则存在整数 \(N\)，使得对每个 \(k\geq1\)，\(E(kN)\) 有一个全局截面，在每个 \(x_i\) 处都是一组基。

可取整数 \(n_0\) 及截面 \(f\in\Gamma(X,L^{\otimes n_0})\)，使每个 \(x_i\) 属于 \(X_f\)，且 \(X_f\) 仿射（EGA II 4.5.4）。在 \(X_f\) 上，由引理 6.7 可取 \(E\) 的截面 \(t\)，在各 \(x_i\) 处均为一组基。又存在整数 \(n_1\)，使截面
\[
 t\otimes f^{\otimes n_1}\in\Gamma(X_f,E(n_0n_1))
\]
能延拓到整个 \(X\)。令 \(N=n_0n_1\)，则 \(f^{\otimes n_1}\in\Gamma(X,L^{\otimes N})\)。将上面所得的 \(E(N)\) 的截面与 \(f^{\otimes kn_1}\) 张量，对每个 \(k\geq0\) 得到 \(E((k+1)N)\) 的截面。由于对每个 \(p\)，有 \(X_{f^{\otimes p}}=X_f\)，而 \(t\) 在每个 \(x_i\) 处为 \(E_{x_i}\) 的一组基，如此构造的 \(E(kN)\) 的截面，在每个 \(x_i\) 处也为一组基。

\sgaStableTarget{sga6.e06.thm.6-9}\textbf{定理 6.9.} 设 \(X\) 是维数为 \(d\) 的诺特概形，且有丰沛层。则
\[
  \Fil^{d+1}K^\bullet(X)=0,
\]
因此 \((\Fil^1K^\bullet(X))^{d+1}=0\)。

（在这里所考虑的一般情形下，后一个等式归功于 J.-P. Serre。）

先证明：若 \(E,E'\) 是 \(X\) 上的两个可逆层，其类仍记为 \(E,E'\)，则 \(E-E'\in K^\bullet(X)^1\)。

设 \(Y\) 是 \(X\) 的一个闭集，\(Y_1,\ldots,Y_k\) 是其不可约分支；对每个 \(i\)，取
\(x_i\in Y_i\setminus\bigcup_{j\ne i}Y_j\)。由引理 6.8，存在整数 \(n,n'\)，使对任意 \(k,k'\geq1\)，可找到从 \(\mathcal O_X\) 到 \(E(kn),E'(k'n')\) 的同态，且在各 \(x_i\) 处均为同构。因此可取态射
\[
 s:L^{\otimes(-nn')}\longrightarrow E,
 \qquad s':L^{\otimes(-nn')}\longrightarrow E',
\]
在各 \(x_i\) 处均为同构。考虑复形
\[
 N^\bullet:\quad
 0\longrightarrow L^{\otimes(-nn')}
 \xrightarrow{(s',0)} E'\oplus L^{\otimes(-nn')}
 \xrightarrow{0+s}E\longrightarrow0.
\]
其类为 \(E-E'\)；此外，\(N^\bullet\) 在各 \(x_i\) 处无环。因此 \(\Supp H^*(N^\bullet)\cap Y\) 不含 \(Y\) 的任何不可约分支，故在 \(Y\) 中的余维数至少为 \(1\)。这就证明了 \(E-E'\in K^\bullet(X)^1\)。

现在取 \(x\in \Fil^{d+1}K^\bullet(X)\)。只需处理乘积
\[
  \gamma^{i_1}(E_1-F_1)\cdots\gamma^{i_k}(E_k-F_k),
  \qquad \sum i_j\geq d+1,
\]
其中 \(E_i,F_i\) 由常秩相同的局部自由层表示。拉回到这些层的旗丛的乘积。此拉回在 \(K^\bullet\) 上为单同态，每个类均分裂为可逆层的类之和。对可逆层的类 \(\xi,\eta\)，
\[
  \gamma_t(\xi-\eta)
  =\frac{\gamma_t(\xi-1)}{\gamma_t(\eta-1)}
  =\frac{1+(\xi-1)t}{1+(\eta-1)t},
\]
所以
\[
  \gamma^i(\xi-\eta)=(-1)^{i-1}(\xi-\eta)(\eta-1)^{i-1}\in K^\bullet(X')^i.
\]
因此，\(x\) 的拉回属于 \(K^\bullet(X')^{d+1}\)。下面由此推出 \(f^*(x)=0\)。设 \(r\) 是 \(X'\) 在 \(X\) 上的相对维数。存在 \(K^\bullet(X')\) 在 \(K^\bullet(X)\) 上的一组基中的一个元素 \(x'\)，形如
\[
 x'=\prod_{i=0}^r(\xi_i-1),
\]
其中各 \(\xi_i\) 是可逆层的类。\footnote{译校注：法文原书第 415 页的乘积下标印作 \(i=0\)，但下文将 \(x'\) 置于第 \(r\) 层；按此处论证，乘积应有 \(r\) 个因子。正文保留原印下标。}要证明这一点，按照 3.1 的论证归结到 \(X'\) 是 \(X\) 上的一个射影丛的情形，再直接应用 1.13 即可。由前面的结论，此 \(x'\) 属于 \(K^\bullet(X')^r\)，因而
\[
 x'f^*(x)\in K^\bullet(X')^{d+r+1}.
\]
但 \(\dim X'=d+r\)，所以由 6.6，\(x'f^*(x)=0\)。由于 \(x'\) 属于 \(K^\bullet(X')\) 在 \(K^\bullet(X)\) 上的一组基，这就推出 \(f^*(x)=0\)。最后，\(f^*\) 由 3.1 为单同态，故 \(x=0\)。

\sgaStableTarget{sga6.e06.problem.6-10}\textbf{问题 6.10.} 在\hyperlink{sga6.e06.thm.6-9}{定理 6.9} 的假设下，是否对每个 \(k\) 都有
\[
  \Fil^kK^\bullet(X)\subset K^\bullet(X)^k?
\]

\sgaStableTarget{sga6.e06.bibliography-for-expose-vi}\subsection*{第 VI 讲参考文献}
\begin{enumerate}[label={\textup{[\arabic*]}}]
\item H. Bass, \emph{K-Theory and stable algebra}, Publications Mathématiques de l'IHÉS, no. 22.
\item C. Chevalley, Séminaire 1958: \emph{Anneaux de Chow et applications}.
\item M. Hakim, \emph{Topos localement annelés et schémas relatifs}, 正在撰写的学位论文。
\item R. Hartshorne, \emph{Residues and duality}, Lecture Notes in Mathematics 20, Springer-Verlag.
\item J.-P. Serre, \emph{Le groupe de Grothendieck des schémas en groupes réductifs}, Publications Mathématiques de l'IHÉS, no. 34.
\end{enumerate}



\providecommand{\Kosz}{\mathrm K}
\providecommand{\Proj}{\operatorname{Proj}}
\providecommand{\Sym}{\operatorname{Sym}}
\providecommand{\gr}{\operatorname{gr}}
\providecommand{\Tor}{\operatorname{Tor}}
\providecommand{\Spec}{\operatorname{Spec}}
\providecommand{\Imf}{\operatorname{Im}}
\providecommand{\Ker}{\operatorname{Ker}}
\providecommand{\Id}{\operatorname{Id}}
\providecommand{\rad}{\operatorname{rad}}
\providecommand{\cO}{\mathcal O}
\providecommand{\cS}{\mathcal S}
\providecommand{\cJ}{\mathcal J}
\providecommand{\cN}{\mathcal N}
\providecommand{\bfP}{\mathbf P}
\providecommand{\bfA}{\mathbf A}


